\ifx\danslelivre\undefined
\documentclass[10pt]{../fichiers-configuration-latex/smfart}
\usepackage{ifthen}
\newboolean{english}
\newboolean{francais}
\setboolean{english}{true}
\setboolean{francais}{false}
\input{../fichiers-configuration-latex/paquets}
\selectlanguage{english}
\input{../fichiers-configuration-latex/lemme-proposition-theoreme}
\input{../fichiers-configuration-latex/numerotation}
\input{../fichiers-configuration-latex/macros}

\externaldocument{../00-introduction/introduction}
\externaldocument{../01-quasi-excellents/excellent}
\externaldocument{../02-topologies/topologies}
\externaldocument{../03-approximation/approximation}
\externaldocument{../04-Cohen-Gabber/Cohen-Gabber}
\externaldocument{../05-algebrisation-partielle/algebrisation}
\externaldocument{../06-log-lissite/log}
\externaldocument{../07-uniformisation-locale/uniformisation-faible}
\externaldocument{../08-modification-theorem/modification}
\externaldocument{../09-uniformisation-premiere-a-l/uniformisation-locale-premiere-a-l}
\externaldocument{../11-produits-orientes/produits-orientes}
\externaldocument{../12-descente-cohomologique-orientee/descente-orientee}
\externaldocument{../13-theoreme-de-finitude/finitude}
\externaldocument{../14-fonctions-dimension/fonctions-dimension}
\externaldocument{../15-Lefschetz-affine/Lefschetz-affine}
\externaldocument{../16-classes-de-Chern/gysin}
\externaldocument{../17-complexes-dualisants/dualite}
\externaldocument{../19-exemples-contre-exemples/contre-exemples}
\externaldocument{../20-rigidite/rigidite}
\externaldocument{../21-finitude-non-abelienne/finitude-non-abelienne}

\begin{document}

\begin{center}
X. Gabber's modification theorem (log smooth case)\\
Luc Illusie and Michael Temkin\footnote{The research of M.T. was partially supported by the
European Union Seventh Framework Programme
(FP7/2007-2013) under grant agreement 268182.}
\end{center}

\tableofcontents

\else
\setboolean{english}{true}
\setboolean{francais}{false}
\selectlanguage{english}
\chapter{Gabber's modification theorem (log smooth case)}
\addtocontents{toc}{par Luc Illusie et Michael Temkin}
\begin{center}
Luc Illusie and Michael Temkin\footnote{The research of M.T. was partially supported by the
European Union Seventh Framework Programme
(FP7/2007-2013) under grant agreement 268182.}
\end{center}
\fi

In this chapter we state and prove a variant of the main theorem of VIII
(see \refext{VIII}{8-TH. MODIFICATION}) for schemes $X$ which are log smooth
over a base $S$ with trivial $G$-action. See \ref{10-modif-logsmooth} for
a precise statement. The proof is given in §1 and in the remaining part
of the exposé we deduce refinements of classical theorems of de Jong,
for schemes of finite type over a field or a trait, where the degree of
the alteration is made prime to a prime $\ell$ invertible on the base.
Sections 2 and 3 are independent and contain two different proofs of
such a refinement, so let us outline the methods briefly.

For concreteness, assume that $k$ is a field, $S=\Spec(k)$, and $X$ is a
separated $S$-scheme of finite type. Two methods to construct regular
$l'$-alterations of $X$ are: (1) use a pluri-nodal fibration to construct
a regular $G$-alteration $X'\to X$ and then factor $X'$ by an $l$-Sylow
subgroup of $G$, and (2) construct a regular $l'$-alteration by induction
on $\dim(S)$ so that one factors by an $l$-Sylow subgroup at each step of
the induction. The first approach is presented in \S\ref{10-primesec}.
It is close in spirit to the approach of \cite{Families@de_Jong} and its
strengthening by Gabber-Vidal, see \cite[\S4]{Brauer@Vidal}. The weak point
of this method is that one uses inseparable Galois alterations. In particular,
even when $k$ is perfect, one cannot obtain a separable alteration of $S$.

In order to deal with the latter case, Gabber suggested to try the second
approach, which is close in spirit to the original de Jong's method of
\cite{Alterations@de_Jong}. Note that the results of \cite{Alterations@de_Jong}
do not provide equivariant alterations, which was the main reason to switch to
the method of \cite{Families@de_Jong}. However, the recent work \cite{Temkin@stable}
resolves this issue, and Gabber's idea was to use \cite{Temkin@stable} instead of
\cite{Alterations@de_Jong}. This approach is realized in \S\ref{10-complsec}
(which is due to the second author) and it outperforms the method of \S\ref{10-primesec}
when $k$ is perfect. Moreover, developing this method the second author discovered
Theorem \ref{10-l'theor} that generalizes Gabber's theorems \ref{10-de Jong1} and
\ref{10-de Jong2} to the case of a general base $S$ satisfying a certain resolvability
assumption (see \S\ref{10-lresolvesec}). In addition, if $S$ is of characteristic zero
then the same method allows to use modifications instead of $l'$-alterations, see
Theorem \ref{10-char0theor}. As an application, in Theorem \ref{10-weaksemistable}
we generalize Abramovich-Karu's weak semistable reduction theorem. Finally,
we minimize separatedness assumptions in \S\ref{10-complsec}, and for this we show
in \S\ref{10-sepelimsec} how to weaken the separatedness assumptions in Theorems
\refext{VIII}{8-TH. MODIFICATION} and \ref{10-modif-logsmooth}.

\section{The main theorem}

\begin{theorem}[Gabber]\label{10-modif-logsmooth}
Let $f : X \ra S$ be an equivariant log smooth map between fs log schemes
endowed with an action of a finite group $G$. Assume that :

(i) $G$ acts trivially on $S$ ;

(ii) $X$ and $S$ are noetherian, qe, separated, log regular, and $f$ defines a
map of log regular pairs $(X,Z) \ra (S,W)$ ;

(iii) $G$ acts tamely and generically freely on $X$.

Let $T$ be the complement of the largest open subset of $X$ over which $G$
acts freely. Then there exists an equivariant projective modification $h :
X' \ra X$ such that, if $Z' = h^{-1}(Z \cup T)$, the pair $(X',Z')$ is log
regular, the action of $G$ on $X'$ is very tame, and $(X',Z')$ is log smooth
over $(S,W)$ as well as the quotient $(X'/G,Z'/G)$ when $G$ acts admissibly
on $X$.
\end{theorem}

\begin{remark2}\label{10-remark main th.} (a) In the absence
of the hypothesis (i) it may not be possible to find a modification $h$
satisfying the properties of \ref{10-modif-logsmooth}, as the example at
the end of \refext{VIII}{8-remarks main th.} shows.

(b) By \cite[8.2]{Toric@Kato} the log smoothness of $f$ and the log regularity
of $S$ imply the log regularity of $X$. Conversely, according to Gabber
(private communication), if $X$ is log regular and $f$ is log smooth and
surjective, then $S$ is log regular.

(c) We will deduce Theorem \ref{10-modif-logsmooth} from Theorem 
\refext{VIII}{8-TH. MODIFICATION}. Recall that in the latter theorem we assumed
that $X$ is qe, though Gabber has a subtler argument that works for a general $X$. 
This forces us to assume that $S$ (and hence $X$) is qe in Theorem 
\ref{10-modif-logsmooth}. However, our argument also shows that
once one removes the quasi-excellence assumption from \refext{VIII}{8-TH. MODIFICATION},
one also obtains the analogous strengthening of Theorem \ref{10-modif-logsmooth}.
\end{remark2}

\medskip For the proof of \ref{10-modif-logsmooth} we will use the following
result on the local structure of equivariant log smooth maps.

\begin{proposition}\label{10-local-structure} (Gabber's preparation
lemma). Let $f : X \ra Y$ be an equivariant log smooth map between fine
log schemes endowed with an action of a finite group $G$. Let $x$ be a
geometric point of $X$, with image $y$ in $Y$. Assume that $G$ is the inertia
group at $x$ and is of order invertible on $Y$. Assume furthermore that $G$
acts trivially on $\overline{M}_{x}$ and $\overline{M}_{y}$\footnote{If $M$
is the sheaf of monoids of a log scheme, $\overline{M}$ denotes, as usual,
the quotient $M/\OO^*$.} and we are given an equivariant chart $a
: Y \ra \Spec \,\Lambda[Q]$ at $y$, modeled on some pairing $\chi : G^{ab}
\otimes Q^{gp} \ra \mu = \mu_N(\mathbf{C})$ (in the sense of (\refext{VI}{6-carte équivariante})),
where $Q$ is fine, $\Lambda = \mathbf{Z}[1/N, \mu]$, with $N$ the exponent of
$G$. Then, up to replacing $X$ by an étale equivariant neighborhood of $x$,
there is an equivariant chart $b : X \ra \Spec \,\Lambda[P]$ extending $a$,
such that $Q^{gp} \ra P^{gp}$ is injective, the torsion of its cokernel is
annihilated by an integer invertible on $X$, and the resulting map $b' :
X \ra X' = Y \times_{\Spec \,\Lambda[Q]}\Spec \,\Lambda[P]$ is smooth. Moreover,
up to further shrinking $X$ around $x$, $b'$ lifts to an equivariant étale
map $c : X \ra X' \times_{\Spec \,\Lambda} \Spec \, \mathrm{Sym}_{\Lambda}(V)$,
where $V$ is a finitely generated projective $\Lambda$-module equipped with a
$G$-action. If $X$, $Y$, and $Q$ are fs, with $Q$ sharp, then $P$ can be chosen
to be fs with its subgroup of units $P^*$ torsionfree.
\end{proposition}

\begin{proof}[Proof of \ref{10-local-structure}.] This is an adaptation of the proof
of \cite[3.5]{LSFI@Kato} to the equivariant case. Consider the canonical
homomorphism of \textit{loc. cit.} $$ k(x) \otimes_{\OO_{X,x}}
\Omega^1_{X/Y,x} \ra k(x) \otimes_{\mathbf{Z}} \overline{M}^{gp}_{X/Y,x}
\leqno(1) $$ sending $1 \otimes d\mathrm{log}t$ to the class of $1 \otimes
t$, where $$ \overline{M}^{gp}_{X/Y,x} = M^{gp}_{X,x}/(\OO_{X,x}^*
+ \mathrm{Im} f^{-1}(M^{gp}_{Y,y})).  $$ It is surjective, and as $G$ fixes
$x$, it is $G$-equivariant. As $G$ is of exponent invertible in $k(x)$
and acts trivially on the right hand side, (1) admits a $G$-equivariant
decomposition $$ k(x) \otimes_{\OO_{X,x}} \Omega^1_{X/Y,x} = V_0
\oplus( k(x) \otimes_{\mathbf{Z}} \overline{M}^{gp}_{X/Y,x}), \leqno(2)
$$ where $V_0$ is a finite dimensional $k(x)$-vector space, endowed with an
action of $G$. Let $(t_i)_{1 \le i \le r}$ be elements of $M^{gp}_x$ such that
the classes of $1 \otimes t_i$ form a basis of $k(x) \otimes_{\mathbf{Z}}
\overline{M}^{gp}_{X/Y,x}$. Let $Z$ be the free abelian group with basis
$(e_i)_{1 \le i \le r}$, and $h : Z \ra M^{gp}_x$ the homomorphism sending
$e_i$ to $t_i$. As $G$ acts trivially on $\overline{M}^{gp}_x$, by the
method of (\refext{VI}{6-preuve très modéré ponctuel}) we get characters
$$ \psi_i : G \ra \mu $$ such that $gh(e_i) = \psi_i(g)h(e_i)$. As in
the proof of \cite[3.5]{LSFI@Kato}, consider the homomorphism $$ u : Z
\oplus Q^{gp} \ra M^{gp}_x $$ defined by $h$ on $Z$ and the composition
$Q^{gp} \ra M^{gp}_y \ra M^{gp}_x$ on the second factor. We have $$ gu(a)
= \psi(g \otimes a))u(a) $$ for some homomorphism $$ \psi : G^{ab} \otimes
(Z \oplus Q^{gp}) \ra \mu $$ extending $\chi$ and such that $\psi(g \otimes
e_i))u(e_i) = \psi_i(g)h(e_i)$. As in \textit{loc. cit.}, if $\overline{u}$
denotes the composition $$ \overline{u} : Z \oplus Q^{gp} \ra M^{gp}_x \ra
\overline{M}^{gp}_x (= M^{gp}_x/\OO^*_x) $$ we see that $k(x) \otimes
\overline{u}$ is surjective, hence the cokernel $C$ of $\overline{u}$ is killed
by an integer $m$ invertible in $k(x)$. Using that $\OO_{X,x}^*$
is $m$-divisible, one can therefore choose elements $a_i  \in M^{gp}_x$ and
$b_i \in Z \oplus Q^{gp}$ ($1 \le i \le n$) such that the images of $a_i$
generate $\overline{M}^{gp}_x$ and $a_i^m = u(b_i)$. Let $E$ be the free
abelian group with basis $e_i$ ($1 \le i \le n$), and let $F$ be the abelian
group defined by the push-out diagram $$ \xymatrix{E \ar[r]^m \ar[d] & E \ar[d]
\\ Z \oplus Q^{gp} \ar[r]^w &F}, \leqno (3) $$ where the left vertical arrow
sends $e_i$ to $b_i$. The lower horizontal map is injective and its cokernel
is killed by $m$. The relation $a_i^m = u(b_i)$ implies that $u$ extends
to a homorphism $$ v : F \ra M^{gp}_x $$ whose composition $\overline{v} :
F \ra M^{gp}_x \ra \overline{M}^{gp}_x$ is surjective. Associated with $v$
is a morphism $$ \varphi : G^{ab} \otimes F \ra \mu $$ extending $\psi$, such
that $gv(a) = \varphi(g \otimes a)v(a)$ for $a \in F$. Let $P := v^{-1}(M_x)
\subset F$. Then $P$ is a fine monoid containing $Q$, $P^{gp} = F$, and $v$
sends $P$ to $M_x$. By \cite[2.10]{LSFI@Kato}, using the ($G$-equivariant)
isomorphism $\mathscr{H}om(\Lambda[P],M_X)_{x} \iso Hom(\Lambda[P],M_{X,x})$,
up to shrinking $X$ at $x$, the equivariant map $$ \Lambda[P] \ra M_x $$
defined by $(v,\varphi)$ extends to an equivariant chart $$ b : X \ra
\Spec \,\Lambda[P] $$ extending the chart $a : Y \ra \Spec \,\Lambda[Q]$. The
homomorphism $Q^{gp} \ra P^{gp}$ is injective, and the torsion part of its
cokernel, which is the same as the torsion part of the cokernel of $w :
Z \oplus Q^{gp} \ra F$ in (3), is killed by $m$. Consider the resulting map
$$ b' : X \ra X' = Y \times_{\Spec \,\Lambda[Q]} \Spec \, \Lambda[P].  $$ This
map is strict. Showing that the underlying schematic map is smooth at $x$ is
equivalent to showing that $b'$ is log smooth at $x$. To do this, as $X$ and
$X'$ are log smooth over $Y$, by the jacobian criterion \cite[3.12]{LSFI@Kato}
it suffices to show that the map $$ k(x) \otimes \Omega^1_{X'/Y} \ra k(x)
\otimes \Omega^1_{X/Y} $$ induced by $b'$ is injective. We have $$ k(x)
\otimes \Omega^1_{X'/Y} = k(x) \otimes P^{gp}/Q^{gp} = k(x) \otimes Z $$
(the last equality by the fact that $F/(Z \oplus Q^{gp})$ is killed by $m$),
and by construction (cf. (2)), we have $$ k(x) \otimes Z = k(x) \otimes
\overline{M}^{gp}_{X/Y,x}, $$ which by the map induced by $b'$ injects into
$k(x) \otimes \Omega^1_{X/Y}$.

Let us now prove the second assertion. For this, as $b'$ is strict, we may
forget the log structures of $X$ and $X'$, and by changing notations, we may
assume that $X' = Y$ and the log structures of $X$ and $Y$ are trivial. In
particular, we have $$ k(x) \otimes \Omega^1_{X/Y} = V_0, $$ with the notation
of (2). As the question is étale local on $X$, and closed points are very
dense in the fiber $X_y$, in particular, any point has a specialization at
a closed point of $X_y$, we may assume that $x$ is localized at a closed
point of $X_y$, and even, up to base changing $Y$ by a finite radicial
extension, that $x$ is a rational point of $X_y$. We then have $$ k(x) \otimes
\Omega^1_{X/Y} = \mathbf{m}_x/(\mathbf{m}_x^2 + \mathbf{m}_y\OO_x),
\leqno (4) $$ where $\mathbf{m}$ denotes a maximal ideal. By \cite[12 2.3,
13 14.4]{Representations@Serre} there is a finitely generated projective
$\Lambda[G]$-module $V$ such that $V_0 = k(x) \otimes V$. The homomorphism $V
\ra \mathbf{m}_x/(\mathbf{m}_x^2 + \mathbf{m}_y\OO_x)$ therefore lifts
to a homomorphism of $\Lambda[G]$-modules $$ V \ra \mathbf{m}_x, \leqno (5) $$
inducing an isomorphism $k(x) \otimes V \ra k(x) \otimes \Omega^1_{X/Y}$. By
the jacobian criterion, it follows that the ($G$-equivariant) map $$ X \ra
Y \times_{\Spec \, \Lambda} \Spec \, Sym_{\Lambda}(V) $$ is étale at $x$.

Let us prove the last assertion. First of all, as $M_x$ is fs, $v : P \ra
M_x$ factors through the saturation $P^{\mathrm{sat}}$ of $P$ in $P^{gp}$,
so we may assume that $P$ is fs. Then (cf.\cite[5.3.42]{HT@GR}) we have
a split exact sequence $$ 0 \ra H \ra P \ra P_0 \ra 0 $$ with $P_0^*$
torsionfree and $H$ a finite group. As $Q$ is fs and sharp, $Q^{gp}$ is
torsionfree, so the composition $Q^{gp} \ra P^{gp} \ra P_0^{gp}$ is still
injective, as well as the composition $H \ra P^{gp} \ra (P^{gp}/Q^{gp})$,
hence $H$ is contained in the torsion part of $(P^{gp}/Q^{gp})$, and
we have an exact sequence $$ 0 \ra H \ra (P^{gp}/Q^{gp})_{tors} \ra
(P_0^{gp}/Q^{gp})_{tors} \ra 0, $$ where the subscript \textit{tors} denotes
the torsion part. Thus $(P_0^{gp}/Q^{gp})_{tors}$ is killed by an integer
invertible on $X$. As $\overline{M}_x$ is torsionfree, the composition
$P \ra M_x \ra \overline{M}_x$ factors through $P_0$, into a map $v_0 :
P \ra \overline{M}_x$. Consider the diagram $$ \xymatrix{{} & M_x \ar[d]
\ar[r] & M^{gp}_x \ar[d] \\ P_0 \ar[r]^{v_0} & \overline{M}_x \ar[r] &
\overline{M}^{gp}_x}, $$ where the square is cartesian. As $P_0^{gp}$
is torsionfree, the map $P^{gp}_0 \ra \overline{M}^{gp}_x$ defined by the
lower row admits a lifting $s : P_0^{gp} \ra M^{gp}_x$, sending $P_0$ to
$M_x$. As $v$ is a chart, $P/v^{-1}(\OO_x^*)  \ra \overline{M}_x$
is an isomorphism, and since $H$ is contained in $v^{-1}(\OO_x^*)$,
$P_0/s^{-1}(\OO_x^*) \ra \overline{M}_x$ is an isomorphism as well,
hence $s$ is a chart at $x$. Associated with $s$ is a homomorphism $$
\varphi_0 : G^{ab} \otimes P_0^{gp} \ra \mu $$ and the map $$ b_0 : X \ra
\Spec \, \Lambda[P_0] $$ defined by the pair $(s,\varphi_0)$ is an equivariant
chart of $X$ at $x$ (extending $a$).
\end{proof}

\medskip \textit{Proof of 1.1 (beginning)}.

\medskip The strategy is to check that, at each step of the proof of the
absolute modification theorem (\refext{VIII}{8-TH. MODIFICATION}), the log
smoothness of $X/S$ is preserved, and, at the end, that of the quotient
$(X/G)/S$ as well. For some of them, this is trivial, as the modifications
performed are log blow ups. Others require a closer inspection.

\subsection{}\label{10-preliminary} \textit{Preliminary reductions.} We
may assume that conditions (1) and (2) at the beginning of  (\refext{VIII}{8-prelimsec})
are satisfied, namely :

(1)\textit{$X$ is regular,}

(2) \textit{$Z$ is a $G$-strict snc divisor in $X$.}

Indeed, these conditions are achieved by $G$-equivariant saturated log blow
up towers (\refext{VIII}{8-step1}, \refext{VIII}{8-step2}).

We will now exploit Gabber's preparation lemma 1.3 to give a local picture
of $f$ displaying both the log stratification and the inertia stratification
of $X$. We work étale locally at a geometric point $x$ in $X$ with image $s$
in $S$. Up to replacing $G$ by the inertia group $G_x$ at $x$, we may assume
that $G = G_x$.

We now apply 1.3. Let $N$ be the exponent of $G$. Assume $S$ strictly
local at $s$. We may replace $\Lambda = \mathbf{Z}[1/N,\mu]$ by its
localization at the (Zariski) image of $s$, so that $\Lambda$ is either the
cyclotomic field $\mathbf{Q}(\mu)$ or its localization at a finite place
of its ring of integers, of characteristic $ p = \mathrm{char}(k(s))$
not dividing $n$. Choose a chart $$ a : S \ra \Spec \,\Lambda[Q] $$ with
$Q$ fs and the inverse image of $\OO_{S,s}^*$ in $Q$ equal
to $\{1\}$, so that $Q$ is sharp and $Q \iso \overline{M}_s$. Let $C$
denote $k(s)$ if $\OO_{S,s}$ contains a field, and a Cohen ring of
$k(s)$ otherwise. Let $(y_i)_{1 \le i \le m}$ be a family of elements of
$\mathbf{m}_s$ such that the images of the $y_i$'s in $\OO_{S,s}/I_s$
form a regular system of parameters, where $I_s = I(s,M_s)$ is the ideal
generated by the image of $M_s - \OO_{S,s}^*$ by the canonical
map $\alpha : M_s \ra \OO_{S,s}$. By \cite[3.2]{Toric@Kato}, the
chart $a$ extends to an isomorphism
\begin{equation}\label{10-Kato-Cohen}
C[[y_1,\cdots,y_m]][[Q]])/(g) \iso \widehat{\OO}_{S,s},
\end{equation}
where $g \in C[[y_1,\cdots,y_m]][[Q]]$ is $0$ if $k(s)$ contains a field,
and congruent to $p = \mathrm{char}(k(s)) > 0$ modulo the ideal generated
by $Q - \{1\}$ and $(y_1,\cdots,y_m)$ otherwise. By 1.3, up to shrinking $X$
around $x$, we can find a $G$-equivariant commutative diagram (with trivial
action of $G$ on the bottom row)
\begin{equation}\label{10-chart-diagram}
\xymatrix{X \ar[dr] \ar[r]^{c} & X' \ar[d] \ar[r]^<<<{b} & \Spec \,(\Lambda[P]
\otimes_{\Lambda}\mathrm{Sym}_{\Lambda}(V))\ar[d] \\ {} & S \ar[r]^a &
\Spec \,\Lambda[Q]},
\end{equation} where :

(i) the square is cartesian ;

(ii) $a$, $b$, and $c$ are strict, where the log
structure on $\Spec \,\Lambda[Q]$ (resp. $\Spec \,(\Lambda[P]
\otimes_{\Lambda}\mathrm{Sym}_{\Lambda}(V))$) is the canonical one, given by
$Q$ (resp. $P$) ; $P$ is an fs monoid, with $P^*$ torsionfree ; $G$ acts on
$\Lambda[P]$ by $g(\lambda p) = \lambda\chi(g,p)p$, for some homomorphism $$
\chi : G^{ab} \otimes P^{gp} \ra \mu $$

(iii) $V$ is a free, finitely generated $\Lambda$-module, equipped with a
$G$-action ;

(iv) the right vertical arrow is the composition of the projection onto the
factor $\Spec \, \Lambda[P]$ and $\Spec \, \Lambda[h]$, for a homomorphisme
$h : Q \ra P$ such that $h^{gp}$ is injective and the torsion part of
$\mathrm{Coker}h^{gp}$ is annihilated by an integer invertible on $X$ ;

(v) $c$ is étale.

\medskip (vi) Consider the map $$ v : P \ra M_x $$ defined by the chart
$X \ra \Spec \, \Lambda[P]$ induced by $bc$. Up to localizing on $X'$
around $x$, we may assume that $v$ factors through the localization
$P_{(\frak p)}$ of $P$ at the prime ideal $\frak p$ complementary of the
face $v^{-1}(\OO_{X,x}^*)$. Replacing $P$ by $P_{(\frak p)}$,
$P$ decomposes into
\begin{equation}\label{10-P1} P = P^* \oplus P_1,
\end{equation} with $P^* = v^{-1}(\OO_{X,x}^*)$ free finitely generated
over $\mathbf{Z}$, and $P_1$ sharp, and the image of $x$ by $bc$ into the
factor $\Spec \, \Lambda[P]$ is the rational point at the origin. Then $v$
induces an isomorphism $P_1 \iso \overline{M}_x$. By the assumptions (1),
(2), we have $\overline{M}_x \iso \mathbf{N}^r$. One can therefore choose
$(e_i \in P_1)$ $(1 \le i \le r)$ forming a basis of $P_1$.  Then $v(e_i) =
t_i \in M_x \subset \OO_{X,x}$ is a local equation for a branch $Z_i$
of $Z$ at $x$, $(Z_i)_{1 \le i \le r}$ is the set of branches of $Z$ at $x$,
and $G$ acts on $t_i$ through the character $\chi_i = \chi | \mathbf{Z}e_i :
G \ra \mu$.

\medskip Furthermore :

\medskip (vii) The square in (\ref{10-chart-diagram}) is tor-independent.

\medskip Indeed, by the log regularity of $S$ and the
choice of the chart $a$, we have, by \cite[6.1]{Toric@Kato},
$Tor_i^{\mathbf{Z[Q]}}(\OO_{S,s},\mathbf{Z}[P]) = 0$ for $i > 0$.

\medskip Though this will not be needed, one can describe the
local structure of (\ref{10-chart-diagram}) more precisely as
follows. Let
\begin{equation}\label{10-Y} Y := \Spec \,(\Lambda[P]
\otimes_{\Lambda}\mathrm{Sym}_{\Lambda}(V)) = \Spec \,(\Lambda[P^*]
\otimes_{\Lambda} \Lambda[P_1]  \otimes_{\Lambda}\mathrm{Sym}_{\Lambda}(V))
\end{equation} and let $Y' := \Spec \,C[[y_1,\cdots, y_m]][[Q]] \times_{\Spec
\Lambda} Y$, with the notation of \ref{10-Kato-Cohen}. Replacing $X'$
by its strict localization at $x$ we may assume that $X = X'$. Then the
completion of $X$ at $x$ is either isomorphic to the completion of $Y'$ at
$x$, or a regular divisor in it, defined by the equation $g' = 0$, where $g'$
is the image of  $g$ in $\widehat{\OO}_{Y',x}$, with the notation
of \ref{10-Kato-Cohen}.

\subsection{}\label{10-Step 3} \textit{Step 3 and log smoothness
(beginning)}. We will now analyze the modifications performed in the
proof of Step 3 in \refext{VIII}{8-step3}, \refext{VIII}{8-step3sec}. The permissible towers used in
\textit{loc. cit.} are iterations of operations of the form : for a subgroup
$H$ of $G$, blow up the fixed point (regular) subscheme $X^H$, and replace $Z$
by the union of its strict transform $Z^{\mathrm{st}}$ and the exceptional
divisor $E$. Though such a blow up is not a log blow up in general, we will
see that it still preserves the log smoothness of $X$ over $S$.

We work étale locally around $x$, so we can assume $X =
X'$ in \ref{10-chart-diagram}. We then have a cartesian square
\begin{equation}\label{10-XH} \xymatrix{X^H \ar[r]^{b^H} \ar[d]^{f} &
Y^H \ar[d] \\ X \ar[r]^b & Y},
\end{equation}
with $Y$ as in (\ref{10-Y}). We
also have cartesian squares
\begin{equation}\label{10-Z} \xymatrix{Z \ar[r]^b
\ar[d]^f & T \ar[d] \\ X \ar[r]^{b} & Y},
\end{equation}
where $T \subset Y$
is the snc divisor $\sum T_i$, $T_i $ defined by the equation $e_i \in P_1$
(\ref{10-P1}), and
\begin{equation}\label{10-Z inter XH} \xymatrix{Z
\times_X X^H \ar[r] \ar[d] & T \times_Y Y^H \ar[d] \\ X \ar[r] & Y}.
\end{equation}
\begin{lemma}\label{10-tor-ind} The squares (\ref{10-XH}),
(\ref{10-Z}), and (\ref{10-Z inter XH}) are tor-independent.
\end{lemma}
\textit{Proof}. For (\ref{10-Z}), this is because $Z$ (resp. $T$) is a
divisor with normal crossings in $X$ (resp. $Y$) (cf. \sga{6}{vii}{1.2}). For
(\ref{10-XH}), as the square (\ref{10-chart-diagram}) is tor-independent
(by \ref{10-preliminary} (vii)), it is enough to show that the composite
(cartesian) square
\begin{equation}\label{10-composite} \xymatrix{X^H
\ar[d] \ar[r] & Y^H \ar[d] \\ S \ar[r] &\Spec \, \Lambda[Q]}
\end{equation}
is tor-independent. We have a decomposition
\begin{equation}\label{10-YH}
Y^H = (\Spec \,\Lambda[P^*]) ^H \times (\Spec \, \Lambda[P_1])^H \times (\Spec \,
\mathrm{Sym}_{\Lambda}(V))^H,
\end{equation}
(products taken over $\Spec \,
\Lambda$), and the map to $\Spec \, \Lambda[Q]$ is the composition of the
projection onto $(\Spec \,\Lambda[P^*]) ^H \times (\Spec \, \Lambda[P_1])^H$
and the canonical map induced by $\Spec \, \Lambda[Q] \ra \Spec \, \Lambda[P]$,
which factors through the fixed points of $H$, $G$ acting trivially on the
base. Let us examine the three factors.

(a) We have $$ (\Spec \,\mathrm{Sym}_{\Lambda}(V))^H = \Spec \,
\mathrm{Sym}_{\Lambda}(V_H), $$ where $V_H$ is the module of coinvariants,
a free module of finite type over $\Lambda$, as $H$ is of order invertible
in $\Lambda$. Therefore $\Spec \, \Lambda[Q] \times_{\Spec \,\Lambda} (\Spec
\,\mathrm{Sym}_{\Lambda}(V))^H$ is flat over $\Spec \, \Lambda[Q]$, and its
enough to check that $(\Spec \,\Lambda[P^*]) ^H \times (\Spec \, \Lambda[P_1])^H$
is tor-independent of $S$ over $\Spec \, \Lambda[Q]$.

(b) The restriction to $P^* = v^{-1}(\OO_{X,x}^*)$ of the
1-cocycle $z(v) \in Z^1(H,Hom(P,k(x)^*)$ associated with $v : P \ra M_x$
($hv(a) = z(v)(h,a)v(a)$ for $h \in H$, $a \in P$, see the proof of
\ref{10-local-structure} and (\refext{VI}{6-preuve très modéré ponctuel})) is a 1-coboundary, hence trivial,
as $B^1(H, Hom(P,k(x)^*) = 0$. Therefore $$
 (\Spec \,\Lambda[P^*]) ^H =  \Spec \,\Lambda[P^*].
$$

(c) Recall that $$ P_1 = \oplus_{1 \le i \le r} \mathbf{N}e_i, $$ with $e_i$
sent by $v$ to a local equation of the branch $Z_i$ of $Z$, and that $G$ acts
on $\Lambda[\mathbf{N}e_i]$ through the character $\chi_i : G \ra \mu$. Let
$A \subset \{1,\cdots,r\}$ be the set of indices $i$ such that $\chi_i | H$
is trivial. Then $$ (\Spec \, \Lambda[P_1])^H = \Spec \, \Lambda[\oplus_{i \in
A} \mathbf{N}e_i].  $$

Let $I$ be the ideal of $P$ generated by $\{e_i \}_{i \notin A}$. It follows
from (b) and (c) that $$ (\Spec \, \Lambda[P])^H = \Spec \, \Lambda[P]/(I),
$$ where $(I)$ is the ideal of $\Lambda[P]$ generated by $I$. By
\cite[6.1]{Toric@Kato}, $Tor_i^{\Lambda[Q]}(\OO_S,\Lambda[P]/(I)) =
0$ for $i > 0$, and therefore (\ref{10-composite}), hence (\ref{10-XH})
is tor-independent. It remains to show the tor-independence of (\ref{10-Z inter XH}).
For this, again it is enough to show the tor-independence
of
\begin{equation}\label{10-new composite} \xymatrix{Z \times_X X^H
\ar[r] \ar[d] & T \times_Y Y^H \ar[d] \\ S \ar[r] & \Spec \, \Lambda[Q]}.
\end{equation} By (a), (b), (c), we have $$ T \times_Y Y^H = \sum_{i \in A}
\Spec \, \Lambda[P]/(J_i) \times \Spec \, \mathrm{Sym}_{\Lambda}(V_H), $$ where
$J_i \subset P$ is the ideal generated by $e_i \in P_1$, and $(J_i)$ the ideal
generated by $J_i$ in $\Lambda[P]$.  The desired tor-independence follows
from the vanishing of $Tor_i^{\Lambda[Q]}(\OO_S,\Lambda[P]/(J_B))$,
where for a subset $B$ of $A$, $J_B$ denotes the ideal generated by the
$e_i$'s for $i \in B$.

\begin{lemma}\label{10-pull-back blow-up} Consider a cartesian square
\begin{equation}\label{10-pull-back 1} \xymatrix{V' \ar[d] \ar[r] & V \ar[d]
\\ X' \ar[r]^g & X},
\end{equation}
where the right vertical arrow is a
regular immersion. If (\ref{10-pull-back 1}) is tor-independent, then the
left vertical arrow is a regular immersion, and $$ \mathrm{Bl}_{V'}(X') = X'
\times_X \mathrm{Bl}_V(X).  $$ Let $W \ra X$ be a second regular immersion,
such that $V \times_X W \ra W$ is a regular immersion, and let $W' = X'
\times_X W$. If moreover the squares
\begin{equation}\label{10-pull-back 2}
\xymatrix{V' \ar[d] \ar[r] & V \ar[d] \\ X' \ar[r]^g & X},
\end{equation}
and
\begin{equation}\label{10-pull-back 3} \xymatrix{V' \times_{X'} W'
\ar[d] \ar[r] & V \times_X W \ar[d] \\ X' \ar[r]^g & X},
\end{equation}
are
tor-independent, then the left vertical arrows are regular immersions, and $$
W'^{\mathrm{st}} = X' \times_X W^{\mathrm{st}}, $$ where $W^{\mathrm{st}}$
(resp. $W'^{\mathrm{st}}$) is the strict transform of $W$ (resp. $W'$)
in $\mathrm{Bl}_V(X)$ (resp. $\mathrm{Bl}_{V'}(X')$).
\end{lemma}
\textit{Proof.} Let $I$ (resp. $I'$) be the ideal of $V$ (resp. $V'$) in $X$
(resp. $X'$). By the tor-independence of (\ref{10-pull-back 1}), if $u :
E \ra I$ is a local surjective regular homomorphism \sga{6}{vii}{1.4}, the
Koszul complex $g^*K(u)$ is a resolution of $\OO_{V'}$, hence $V'
\ra X'$ is a regular immersion. Moreover, by \sga{6}{vii}{1.2}, for any $n
\ge 0 $, the natural map $g^*I^n \ra I'^n$ is an isomorphism,  and therefore
$\mathrm{Bl}_{V'}(X') = X' \times_X \mathrm{Bl}_V(X)$. The tor-independence of
(\ref{10-pull-back 2}) and (\ref{10-pull-back 3}) imply that of $$ \xymatrix{
V' \times_{X'} W' \ar[r] \ar[d] & V \times_X W \ar[d] \\ W' \ar[r] & W}.
$$ The second assertion then follows from the first one and the formulas
(\refext{VIII}{8-strictrem} (ii)) $$ W^{\mathrm{st}} = \mathrm{Bl}_{V \times_X
W} W, $$ $$ W'^{\mathrm{st}} = \mathrm{Bl}_{V' \times_{X'} W'} W'.  $$

\subsection{}\label{10-Step 3 (end)} \textit{Step 3 and log smoothness
(end)}. As recalled at the beginning of \ref{10-Step 3}, we have to show
that, if $H$ is a subgroup of $G$, then the log regular pair $(X_1,Z_1)$
is log smooth over $S$, where $X_1 := \mathrm{Bl}_{X^H}(X)$ and $Z_1$ is
the snc divisor $Z^{\mathrm{st}} \cup E$, $Z^{\mathrm{st}}$ (resp. $E$)
denoting the strict transform of $Z$ (resp. the exceptional divisor) in the
blow-up $h : X_1 \ra X$.

The question is again étale local above $X$ around $x$, so we may assume
that $X = X'$ and we look at the cartesian square (\ref{10-chart-diagram}) $$
\xymatrix{X \ar[r] \ar[d] & Y \ar[d] \\ S \ar[r] & \Spec \, \Lambda[Q]}, $$ with
$Y$ as in (\ref{10-Y}), and the associated cartesian squares (\ref{10-XH}),
(\ref{10-Z}), and (\ref{10-Z inter XH}).

\textit{Claim.} We have
\begin{equation}\label{10-blow-up}
\mathrm{Bl}_{X^H}(X) = X \times_Y \mathrm{Bl}_{Y^H}(Y),
\end{equation}
\begin{equation}\label{10-strict transform} Z^{\mathrm{st}} = X \times_Y
T^{\mathrm{st}}.
\end{equation}

\textit{Proof.} In view of \ref{10-tor-ind} and \ref{10-pull-back blow-up},
(\ref{10-blow-up}) follows from the fact that the immersion $Y^H \ra
Y$ is regular. For (\ref{10-strict transform}) recall that $$ T = T_0
\times_{\Spec \, \Lambda} \Spec \, \mathrm{Sym}_{\Lambda}(V), $$ where $T_0
\subset \Spec \, \Lambda[P]$ is the snc divisor $$ T_0 = \sum_{1 \le i \le r}
\mathrm{div}(z_i) $$ with $z_i \in \Lambda[P]$ the image of $e_i \in P_1$
as in \ref{10-P1}. Hence
\begin{equation}\label{10-T} T = \sum_{1 \le
i \le r} T_i,
\end{equation}
where $T_i = \mathrm{div}(z_i) \times_{\Spec \,
\Lambda}\mathrm{Sym}_{\Lambda}(V)$, and $T^{\mathrm{st}} = \sum_{1 \le i \le
r}T_i^{\mathrm{st}}$. We have (\ref{10-YH}) $$ Y^H =  (\Spec \, \Lambda[P])^H
\times \Spec \, \mathrm{Sym}_{\Lambda}(V_H), $$ with $(\Spec \, \Lambda[P])^H$
defined by the equations $(z_i = 0)_{i \notin A}$, with the notations of
\ref{10-tor-ind} (c). In particular, the immersion $Y^H \times_Y Z_i \ra Z_i$
is regular, hence, by \ref{10-pull-back blow-up}, we have $Z_i^{\mathrm{st}}
= X \times_Y T_i^{\mathrm{st}}$, hence (\ref{10-strict transform}), which
finishes the proof of the claim.


Since the map $S \ra \Spec \, \Lambda[Q]$ is strict, in order to prove the
desired log smoothness, we may, by this claim, replace the triple $(X, X^H,
Z)$ over $S$ by $(Y,Y^H,T)$ over $\Spec \, \Lambda[Q]$. We choose coordinates on
$P^*$, $P_1 = \mathbf{N}^r$, $V$ : $$ P^* = \oplus_{1 \le i \le t}\mathbf{Z}f_i
, \, \, P_1 = \oplus_{1 \le i \le r}\mathbf{N}e_i, \, \, V = \oplus_{1 \le i
\le s} \Lambda y_i $$ $$ \Lambda[P] = \Lambda[u_1^{\pm 1}, \cdots, u_t^{\pm 1},
z_1, \cdots, z_r], \, \, \mathrm{Sym}_{\Lambda}(V) = \Lambda[y_1,\cdots, y_s],
$$ with $u_i$ (resp. $z_i$) the image of $f_i$ (resp. $e_i$) in $\Lambda[P]$,
in such a way that $$ \Lambda[P]^H = \Lambda[u_1^{\pm 1}, \cdots, u_t^{\pm 1},
z_{m+1}, \cdots, z_r], $$ i. e. is defined in $\Lambda[P]$ by the equations
$(z_1 = \cdots = z_m = 0$, for some $m$,  $1 \le m \le r$, and $$ \Lambda[V_H]
= \Lambda[y_{n+1}, \cdots, y_s], $$ i. e. is defined in $\Lambda[V]$ by the
equations $y_1 = \cdots = y_n = 0$ for some $n$, $1 \le n \le s$. Then $$
Y^H \subset Y = \Spec \, \Lambda[u_1^{\pm 1}, \cdots, u_t^{\pm 1},z_1, \cdots,
z_r, y_1, \cdots, y_s] $$ is defined by the equations $$ z_1 = \cdots =
z_m = y_1 = \cdots = y_n = 0.  $$ Then$$ Y' := \mathrm{Bl}_{Y^H}(Y) $$ is
covered by affine open pieces : $$ U_i = \Spec \, \Lambda[(u_j^{\pm 1})_{1
\le j  \le t}, z'_1,\cdots, z'_{i-1}, z_i, z'_{i+1}, \cdots, z'_m, z_{m+1},
\cdots, z_r, y'_1, \cdots, y'_n, y_{n+1}, \cdots, y_s] $$ ($1 \le i \le m$),
with $U_i \ra Y$ given by $z_j \ra z_iz'_j$ for $1 \le j \le m$, $j \ne i$,
$y_j \ra z_iy'_j$, $1 \le j \le n$, and the other coordinates unchanged,
and $$ V_i = \Spec \, \Lambda[(u_j^{\pm 1})_{1 \le j  \le t}, z'_1, \cdots,
z'_m,  z_{m+1}, \cdots, z_r, y'_1, \cdots, y'_{i-1},y_i,y'_{i+1}, \cdots,
y'_n, y_{n+1}, \cdots, y_s] $$ ($1 \le i \le n$), with $V_i \ra Y$ given by
$z_j \mapsto y_iz'_j$ for $1 \le j \le m$, $y_j \mapsto y_iy'_j$, $1 \le j \le
n$, $j \ne i$, and the other coordinates unchanged.  Recall that $Y$ has the
log structure defined by the log regular pair $(Y,T)$, where $T$ is the snc
divisor $$ T = (z_1 \cdots z_r = 0), $$ and $Y'$ is given the log structure
defined by the log regular pair $(Y',T')$, where $T'$ is the snc divisor $$
T' = F \cup T^{\mathrm{st}}, $$ where $F$ is the exceptional divisor of the
blow up of $Y^H$ and $T^{\mathrm{st}}$ the strict transform of $T$. Consider
the canonical morphisms $$ \xymatrix{Y' \ar[r]^b & Y \ar[r]^<<<{g} & \Sigma :=
\Spec \, \Lambda[Q]}.  $$ They are both morphisms of log schemes. The morphism
$g$ is given by the homomorphism of monoids $\gamma : Q \ra P$, i. e.  $$
q \in Q \mapsto (\gamma_1(q), \cdots, \gamma_t(q), \gamma_{t+1}(q), \cdots,
\gamma_{t+r}(q), 0, \cdots, 0) \in \Lambda[u_1^{\pm 1}, \cdots, u_t^{\pm 1},
z_1, \cdots, z_r, y_1, \cdots, y_s].  $$ The blow up $b$ has been described
above in the various charts. Note that $b$ is not log étale, or even log
smooth, in general. However, the composition $gb : Y' \ra \Sigma$ is log
smooth. We'll check this on the charts $(U_i)$, $(V_i)$.

(a) \textit{Chart of type $U_i$}. We have $F = (z_i = 0)$, $T^{\mathrm{st}}
= (\prod_{1 \le j \le r, j \ne i}z'_i = 0)$. Hence the log strucure of
$U_i$ is given by the canonical log structure of $\Lambda[\mathbf{N}^r]$
in the decomposition $$ U_i = \Spec \, \Lambda[\mathbf{Z}^t] \times \Spec \,
\Lambda[\mathbf{N}^r] \times \Spec \, \Lambda[y'_1, \cdots, y'_n, y_{n+1},
\cdots, y_s] $$ with the basis element $e_k$ of $\mathbf{N}^r$ sent to the
$k$-th place in $(z'_1,\cdots, z'_{i-1}, z_i, z'_{i+1}, \cdots, z'_m, z_{m+1},
\cdots, z_r)$ (and the basis element $f_k$ of $\mathbf{Z}^t$ sent to $u_k$),
the third factor having the trivial log structure. Checking the log smoothness
of $gb : U_i \ra \Sigma$ amounts to checking the log smoothness of its factor
$\Spec \, \Lambda[P] \ra \Sigma = \Spec \, \Lambda(Q)$, which is defined by the
composition of homomorphisms of monoids $$ \xymatrix{Q \ar[r]^{\gamma} &
\mathbf{Z}^t \oplus  \mathbf{N}^r \ar[r]^{Id \oplus \beta} & \mathbf{Z}^t
\oplus \mathbf{N}^r}, $$ where $\beta$ is the homomorphism $\mathbf{N}^r
\ra \mathbf{N}^r$ sending $e_j$ to $e_j + e_i$ for $1 \le j \le m, j
\ne i$, $e_i$ to $e_i$, and $e_j$ to $e_j$ for $m+1 \le j \le r$. Recall
((\ref{10-chart-diagram}), (iv)) that $\gamma^{gp}$ is injective and the
torsion part of its cokernel is invertible in $\Lambda$. As $\beta^{gp}$ is
an isomorphism, the same holds for the composition $(Id \oplus \beta)\gamma$,
hence $gb : U_i \ra \Sigma$ is log smooth.

\medskip (b) \textit{Chart of type $V_i$}. We have $F = (y_i = 0)$,
$T^{\mathrm{st}} = \prod_{1 \le j \le m}z'_j\prod_{j \ge m+1}z_i$. Hence
the log structure of $V_i$ is given by the canonical log structure
of $\Lambda[\mathbf{N}^{r+1}]$ in the decomposition $$ V_i = \Spec \,
\Lambda[\mathbf{Z}^t] \times \Spec \, \Lambda[\mathbf{N}^{r+1}] \times \Spec
\, \Lambda[(y'_j)_{1 \le j \le n, j \ne i}, y_{n+1}, \cdots, y_s] $$
with the basis element $e_k$ of $\mathbf{N}^{r+1}$ sent to the $k$-th
place in $(z'_1,\cdots, z'_m, z_{m+1}, \cdots, z_r)$ if $k \le r$, and
$e_{r+1}$ sent to $y_i$ (and the basis element $f_k$ of $\mathbf{Z}^t$
sent to $u_i$), the third factor having the trivial log structure. Again,
Checking the log smoothness of $gb : V_i \ra \Sigma$ amounts to checking
the log smoothness of its factor $\Spec \, \Lambda[\mathbf{Z}^t] \times \Spec \,
\Lambda[\mathbf{N}^{r+1}] \ra \Spec \, \Lambda(Q)$. This factor is defined by
the composition of homomorphisms of monoids $$ \xymatrix{Q \ar[r]^{\gamma}
& \mathbf{Z}^t \oplus \mathbf{N}^r \ar[r]^{Id \oplus \beta} & \mathbf{Z}^t
\oplus \mathbf{N}^{r+1}} $$ where $\beta : \mathbf{N}^r \ra \mathbf{N}^{r+1}$
sends $e_j$ to $e_j + e_{r+1}$ for $1 \le j \le m$, and to $e_j$ for  $m+1
\le j \le r$. Then $\beta^{gp}$ is injective, and its cokernel is isomorphic
to $\mathbf{Z}$, hence $(\beta \gamma)^{gp}$ is injective, and we have an
exact sequence $$ 0 \ra \mathrm{Coker}\,  \gamma^{gp} \ra \mathrm{Coker} \,
(\beta \gamma)^{gp} \ra \mathbf{Z} \ra 0.  $$ In particular, the torsion
part of $\mathrm{Coker} (\beta \gamma)^{gp}$ is isomorphic to that of
$\mathrm{Coker} \, \gamma^{gp}$, hence of order invertible in $\Lambda$,
which implies that $gb : V_i \ra \Sigma$ is log smooth.

This finishes the proof that Step 3 preserves log smoothness.

\subsection{}\label{10-end of proof} \textit{End of proof of
\ref{10-modif-logsmooth}.} We may now assume that in addition to conditions
(1) and (2) of \ref{10-preliminary}, condition (3) is satisfied as well, namely

(3) \textit{$G$ acts freely on $X - Z$ (i. e. $Z = Z \cup T$ in the notation
of \ref{10-modif-logsmooth} or (\refext{VIII}{8-TH. MODIFICATION})),  and,
for any geometric point $x \ra X$, the inertia group $G_x$ is abelian.}

We have to check :

\medskip \textit{Claim.} \textit{If $f_{(G,X,Z)} : (X',Z') \ra (X,Z)$ is the
modification of (\refext{VIII}{8-logquotcor}),
then $(X',Z')$ and $(X'/G,Z'/G)$ are log smooth over $S$.}

\medskip Working étale locally around a geometric point $x$ of $X$, we will
first choose a strict rigidification $(X,\overline{Z})$ of $(X,Z)$ such that
$(X,\overline{Z})$ is log smooth over $S$.  We will define $(X,\overline{Z})$
as the pull-back by $S \ra\Sigma = \Spec \, \Lambda[Q]$ of a rigidification
$(Y,\overline{T})$ of $(Y,T)$ which is log smooth over $\Sigma$, with the
notation of $(\ref{10-Y})$. Using that $G$  $(= G_x)$ is abelian, one can
decompose $V$ into a sum of $G$-stable lines, according to the characters
of $G$ : $$ V = \oplus_{1 \le i \le s} \Lambda y_i $$ with $G$ acting on
$\Lambda y_i$ through a character $\chi_i : G \ra \mu_N$, i. e. $gy_i  =
\chi_i(g)y_i$. We define $\overline{T}$ to be the divisor $z_1 \cdots z_r y_1
\cdots y_s = 0$ in $Y = \Spec \, \Lambda[u_1^{\pm 1}, \cdots, u_t^{\pm 1},z_1,
\cdots, z_r, y_1, \cdots, y_s]$. The action of $G$ on $(Y,\overline{T})$ is
very tame at $x$ because the log stratum at $x$ is $\Spec \, \Lambda[u_1^{\pm
1}, \cdots, u_t^{\pm 1}]$, hence very tame in a neighborhood of $x$ by
(\refext{VIII}{8-tamelocallem}) (actually on the whole of $Y$,
cf. (\refext{VIII}{6-dcn stricts}, \refext{VIII}{6-exemples}(a)).
On the other hand, $(\Spec \, \Lambda[y_1,\cdots, y_s], y_1 \cdots y_s
= 0)$ is log smooth over $\Spec \,\Lambda$, and as $\Spec \, \Lambda[P]$ is log
smooth over $\Sigma$, $(Y,\overline{T}) $ is log smooth over $\Sigma$. Since
$f_{(G,X,Z)}$ is compatible with base change by strict inert morphisms,
it is enough to check that if $f_{(G,Y,T)} = f_{(G,Y,T,\overline{T})}
: (Y',T') \ra (Y,T)$ is the modification of (\refext{VIII}{8-logquotcor})
then $(Y',T')$ is log smooth over $\Sigma$. Recall
(\refext{VIII}{8-constrsec}) that we have a cartesian $G$-equivariant diagram
\begin{equation}\label{10-cartesian diagram} \xymatrix{(Y',\overline{T'})
\ar[r]^{\overline{h'}} \ar[d]^{\alpha'} & (Y,\overline{T}) \ar[d]^{\alpha}
\\ (Y'_1,\overline{T'_1}) \ar[r]^{\overline{h_1}} & (Y_1,\overline{T_1})},
\end{equation} where the horizontal maps are the compositions of saturated
log blow up towers, and the vertical ones Kummer étale $G$-covers. From
(\ref{10-cartesian diagram}) is extracted the relevant diagram involving $h
:= f_{(G,Y,T,\overline{T})}$, $$ \xymatrix{(Y',T') \ar[d]^{\beta} \ar[r]^h &
(Y,T) \\ (Y'_1,T'_1)}, $$ where $T'_1 = \overline{h}^{-1}_1(T_1)$, with $T_1 =
T/G$, $T' = \alpha'^{-1}(T'_1)$, and $h$ (resp. $\beta$) is the restriction
of $\overline{h'}$ (resp. $\alpha'$) over $(Y,T)$ (resp. $(Y_1',T_1')$). In
particular, $\beta$ is a Kummer étale $G$-cover (as Kummer étale
$G$-covers are stable under any fs base change). As $G$ acts trivially
on $S$, this diagram can be uniquely completed into a commutative diagram
\begin{equation}\label{10-relevant diagram} \xymatrix{(Y',T') \ar[d]^{\beta}
\ar[r]^h & (Y,T) \ar[d]^f \\ (Y'_1,T'_1) \ar[r]^g & \Sigma}.
\end{equation}
Here $f$ is log smooth and $\beta$ is a Kummer étale $G$-cover. Though
$\overline{h'}$ and $\overline{h_1}$ are log smooth, $h$ and $h_1$ are not,
in general. However, it turns out that :

\medskip (*) \textit{$g : (Y'_1,T'_1) \ra \Sigma$, hence $g\beta = fh :
(Y',T') \ra \Sigma$, are log smooth},

\medskip \noindent which will finish the proof of the claim, hence of
\ref{10-modif-logsmooth}. We first prove

\medskip (**) \textit{With the notation of (\ref{10-relevant diagram}),
$(Y_1,\overline{T_1})$ is log smooth over $\Sigma$}.

\medskip Let us write $Y = \Spec \, \Lambda[\overline{P}]$, with
\begin{equation}\label{10-Peq} \overline{P} = P \times \mathbf{N}^s =
\mathbf{Z}^t \times \mathbf{N}^r \times \mathbf{N}^s.
\end{equation} As
$G$ acts very tamely on $(Y,\overline{T})$, the quotient pair $(Y_1 = Y/G,
\overline{T_1} = \overline{T}/G)$ is log regular. More precisely, by the
calculation in (\refext{VI}{6-très modéré ponctuel}(b)), this pair consists of the log scheme $Y_1 =
\Spec \, \Lambda[\overline{R}]$ with its canonical log structure, where $$
\overline{R} = \mathrm{Ker}(\overline{P}^{gp} \ra \mathrm{Hom}(G,\mu_N))
\cap \overline{P}, $$ $\overline{P}^{gp} \ra \mathrm{Hom}(G,\mu_N)$ being
the homomorphism defined by the pairing $\chi : G \otimes \overline{P}^{gp}
\ra \mu_N$. The inclusion $\overline{R} \subset \overline{P}$ is a Kummer
morphism, and $\overline{P}^{gp}/\overline{R}^{gp}$ is annihilated by
an integer invertible in $\Lambda$. As $Q^{gp} \ra P^{gp}$ is injective,
with the torsion part of its cokernel annihilated by an integer invertible
in $\Lambda$, the same is true of $Q^{gp} \ra \overline{P}^{gp}$, hence
of $Q^{gp} \ra \overline{R}^{gp}$. Thus $(Y_1,\overline{T_1}) = \Spec \,
\Lambda[\overline{R}]$ is log smooth over $\Sigma$.

Finally, let us prove (*). It is enough to work locally on $Y'_1$ so we
can replace the log blow up sequence $(Y'_1,\oT'_1)\to(Y_1,\oT_1)$ with
an affine chart (i.e. we replace the first log blow up with a chart, then
do the same for the second one, etc.). Then $Y'_1=\Spec\Lambda[\oR']$,
and $\oR^\gp\iso\oR'^\gp$ by \refext{VIII}{8-blowmonosec}. Note that
$\oR'\iso\ZZ^a\times\NN^b$ where $D_1,\dots,D_b$ are the components of $\oT'_1$.
We can assume that $D_1,\dots,D_c\subseteq T'_1$ and $D_{c+1},\dots,D_b$
are not contained in $T'_1$. Let $R'\iso\ZZ^a\times\NN^c$ denote the
submonoid $\oR'$ that defines the log structure of $(Y'_1,T'_1)$. Note that
$R'$ consists of all elements $g'\in\oR'$ such that $(g'=0)\subseteq T'_1$
(as a set). Also, by $\nu:\oR\to\oR'$ we will denote the homomorphism that
defines $(Y'_1,\oT'_1)\to(Y_1,\oT_1)$.

We showed in \refext{VIII}{8-constrsec} that $T_1=T/G$ is a $\QQ$-Cartier divisor in $Y_1$
and observed that therefore $T'_1$ is a Cartier divisor in $Y'_1$. Note that
the inclusion $R \subset \overline{R}$, where $$ R = \mathrm{Ker}(P^{gp}
\ra \mathrm{Hom}(G,\mu_N)) \cap \overline{P} $$ defines a log structure
on $Y_1$. Denote the corresponding log scheme $(Y_1,T_1)$. We obtain the
following diagram of log schemes (on the left). The corresponding diagram of
groups is placed on the right; we will use it to establish log smoothness
of $g$. Existence of dashed arrows requires an argument; we will construct
them later.

\begin{equation}\label{10-doublediagram} \xymatrix{(Y'_1,\oT'_1)
\ar[d] \ar[r]^{\oh_1} & (Y_1,\oT_1) \ar[d]\ar[dr] &
& & \oR'^\gp&\oR^\gp\ar[l]_\sim^{\nu^\gp}\\ (Y'_1,T'_1)
\ar@{-->}[r] & (Y_1,T_1)\ar[r]& \Sigma& & R'^\gp\ar@{^{(}->}[u]&
R^\gp\ar@{_{(}-->}[l]\ar@{^{(}->}[u]&Q^\gp\ar[l]\ar[lu]}
\end{equation}

Part (ii) of the following remark clarifies the notation $(Y_1,T_1)$. It
will not be used so we only sketch the argument.

\begin{remark} (i) Note that $(Y_1,T_1)$ may be not log smooth over
$\Sigma$. For example, even when $\Sigma$ is log regular, e.g. $\Spec k$ with
trivial log structure, $(Y_1,T_1)$ does not have to be log regular, as $T_1$
may even be non-Cartier. Nevertheless, as $\overline{h_1}$ is log smooth (even
log étale), $(Y'_1,\overline{T'_1})$ is log smooth over $\Sigma$. Moreover,
$Y'_1$ is regular, and $\overline{T'_1}$ an snc divisor in it.

(ii) Although $T_1$ may be bad, one does have that
$R\calO_{Y_1}^\times=\calO_{Y_1}\cap i_*\calO^\times_{Y_1\setminus
T_1}$ for the embedding $i:Y_1\setminus T_1\into Y_1$. This can
be deduced from the formulas for $R$ and $\oR$ and the fact that
$\oR\calO_{Y_1}^\times=\calO_{Y_1}\cap j_*\calO^\times_{Y_1\setminus\oT_1}$
by log regularity of $(Y_1,\oT_1)$.
\end{remark}

Note that $Q\to\oP$ factors through $P$, hence $Q\to\oR$ factors through
$R=\oR\cap P$. It follows from (\ref{10-Peq}) that $P$ consists of all elements
$f\in\oP$ whose divisor $(f=0)$ is contained in $T$ (as a set). Therefore
$g\in\oR$ lies in $R$ if and only if $(g=0)\subseteq T_1$ (as a set). This
fact and the analogous description of $R'$ observed earlier imply that
$\nu:\oR\to\oR'$ takes $R$ to $R'$. Thus, we have established the dashed
arrows in (\ref{10-doublediagram}).

Let $\varphi:Q\to\oR'$ be the homomorphism defining the composition
$(Y'_1,\overline{T'_1}) \ra (Y_1,\overline{T_1})\to\Sigma$. Since the
latter is log smooth, $\varphi$ is injective, and the torsion part
of $\Coker(\varphi^\gp)$ is annihilated by an integer $m$ invertible
in $\Lambda$.  Note that $R^\gp\into R'^\gp\into\oR'^\gp$, and therefore
we also have that $Q^\gp\into R'^\gp$ and the torsion of its cokernel is
annihilated by $m$. Therefore, $(Y'_1,T'_1)$ is log smooth over $\Sigma$,
which finishes the proof of (*), hence of $\ref{10-modif-logsmooth}$.

\begin{remark}\label{10-remark proof of (*)} In the proof of (*) above, we
first proved that $g$ is log smooth, and deduced that $g\beta$ is, too. In
fact, as $\beta$ is a Kummer étale $G$-cover, the log smoothness of $g\beta$
implies that of $g$. More generally, we have the following descent result,
due to Kato-Nakayama (\cite[3.4]{quasisections@Nakayama}) :
\end{remark}

\begin{theorem} Let $\xymatrix{X' \ar[r]^g & X \ar[r]^f & Y}$ be morphisms
of fs log schemes. If $g$ is surjective, log étale and exact,  and $fg$
is log smooth, then $f$ is log smooth.
\end{theorem} The assumption
on $g$ is equivalent to saying that $g$ is a Kummer étale cover
(cf. \cite[1.6]{Overview@Illusie}).

\section{Prime to $\ell$ variants of de Jong's alteration theorems}\label{10-primesec}

Let $X$ be a noetherian scheme, and $\ell$  be a prime number. A morphism
$h : X' \ra X$ is called an \textit{$\ell'$-alteration} if $h$ is proper,
surjective, generically finite, maximally dominating (i.e., (\refext{II}{2-maximalement-dominant})
sends each maximal point to a maximal point) and the degrees of
the residual extensions $k(x')/k(x)$ over each maximal point $x$ of $X$
are prime to $\ell$. The next theorem was stated in (Introduction 0.3 (1)) :

\begin{theorem}[Gabber]\label{10-de Jong1} Let $k$ be a field, $\ell$ a
prime number different from the characteristic of $k$, $X$ a separated and
finite type $k$-scheme, $Z \subset X$ a nowhere dense closed subset. Then
there exists a finite extension $k'$ of $k$, of degree prime to $\ell$,
and a projective $\ell'$-alteration $h : \tilde X \ra X$ above $\Spec \,
k' \ra \Spec \, k$, with $\tilde X$ smooth and quasi-projective over $k'$,
and $h^{-1}(Z)$ is the support of a relative, strict normal crossings divisor.
\end{theorem} Recall that a relative, strict normal crossings divisor in
a smooth scheme $T/S$ is a divisor $D = \sum_{i \in I}D_i$, where $I$ is
finite, $D_i \subset T$ is an $S$-smooth closed subscheme of codimension 1,
and for every subset $J$ of $I$ the scheme-theoretic intersection $\cap_{i
\in J}D_i$ is smooth over $S$.

\medskip We will need the following variant, due to Gabber-Vidal (proof
of \cite[4.4.1]{BS@Vidal}), of de Jong's alteration theorems \cite[5.7,
5.9, 5.11]{Families@de_Jong}, cf. \cite[3.8]{Independance@Zheng} :
\begin{lemma}\label{10-de Jong-Gabber-Vidal1} Let $X$ be a proper scheme over
$S = \Spec \, k$, normal and geometrically reduced and irreducible, $Z \subset X$
a nowhere dense closed subset. We assume that a finite group $H$ acts on $X \ra
S$, faithfully on $X$, and that $Z$ is $H$-stable. Then there exists a finite
extension $k_1$ of $k$, a finite group $H_1$, a surjective homomorphism $H_1
\ra H$, and an $H_1$-equivariant diagram with a cartesian square (where $S =
\Spec \, k$, $S_1 = \Spec \, k_1$)
\begin{equation}\label{10-Vidal-diagram1}
\xymatrix{X \ar[d] & X_1 \ar[d] \ar[l]_b & X_2 \ar[l]_a \ar[dl] \\ S &
S_1 \ar[l] & {}}
\end{equation} satisfying the following properties :

(i) $S_1/\mathrm{Ker}(H_1 \ra H) \ra S$ is a radicial extension ;

(ii) $X_2$ is projective and smooth over $S_1$ ;

(iii) $a : X_2 \ra X_1$ is projective and surjective, maximally dominating
and generically finite and flat,  and there exists an $H_1$-admissible dense
open subset $W \subset X_2$ over a dense open subset $U$ of $X$, such that
if $U_1 = S_1 \times_S U$ and $K = \mathrm{Ker}(H_1 \ra \mathrm{Aut}(U_1))$,
$W \ra W/K$ is a Galois étale cover of group $K$ and the morphism $W/K \ra
U_1$ induced by $a$ is a universal homeomorphism ;

(iv) $(ba)^{-1}(Z)$ is the support of a strict normal crossings divisor
in $X_2$.
\end{lemma}

\textit{Proof.} We may assume $X$ of dimension $d \ge 1$. We apply
\cite[4.4.3]{BS@Vidal} to $X/S$, $Z$, and $G = H$. We get the data of
(\textit{loc. cit.}), namely an equivariant finite extension of fields
$(S_1, H_1) \ra (S,H)$ such that $S_1/\mathrm{Ker}(H_1 \ra H) \ra S$ is
radicial, an $H_1$-equivariant pluri-nodal fibration $(Y_d \ra \cdots \ra
Y_1 \ra S_1, \{ \sigma_{ij} \}, Z_0 = \emptyset)$, and an $H_1$-equivariant
alteration $a_1 : Y_d \ra X$ over $S$, satisfying the conditions (i),
(ii), (iii) of (\textit{loc. cit.}) (in particular $a_1^{-1}(Z) \subset
Z_d$). Then, as in the proof of \cite[4.4.1]{BS@Vidal}, successively applying
\cite[4.4.4]{BS@Vidal} to each nodal curve $f_i : Y_i \ra Y_{i-1}$, one can
replace $Y_i$ by an $H_1$-equivariant projective modification $Y'_i$ of it
such that $Y'_i$ is regular, and the inverse image $Z'_i$ of $Z_i := \cup_j
\sigma_{ij}(Y_{i-1}) \cup f_i^{-1}(Z_{i-1})$ in $Y'_i$ is an $H_1$-equivariant
strict snc divisor. Then,  $X_2 := Y'_d$ is smooth over $S_1$ and $Z'_d$
is a relative snc divisor over $S_1$. This follows from the analog of
the remark following \cite[4.4.4]{BS@Vidal} with ``semistable pair over
a trait" replaced by ``pair consisting of a smooth scheme and a relative
snc divisor over a field". In particular, $(ba)^{-1}(Z)_{\mathrm{red}}$
is a subdivisor of $Z'_d$, hence an snc divisor. After replacing $H_1$ by
$H_1/\mathrm{Ker}(H_1 \ra \mathrm{Aut}(X_2))$ the open subsets $U$ and $V$
of (iii) are obtained as at the end of the proof of \cite[4.4.1]{BS@Vidal}.

\subsection{}\label{10-proof of de Jong1} \textit{Proof of \ref{10-de
Jong1}}. There are three steps.

\medskip \textit{Step 1}. \textit{Preliminary reductions.} By Nagata's
compactification theorem \cite{Nagata@Conrad}, there exists a dense open
immersion $X \subset \overline{X}$ with $\overline{X}$ proper over $S$. Up
to replacing $X$ by $\overline{X}$ and $Z$ by its closure $\overline{Z}$,
we may assume $X$ proper over $S$. By replacing $X$ by the disjoint sum
of its irreducible components, we may further assume $X$ irreducible, and
geometrically reduced (up to base changing by a finite radicial extension
of $k$). Up to blowing up $Z$ in $X$ me may further assume that $Z$ is a
(Cartier) divisor in $X$. Finally, replacing $X$ by its normalization $X'$,
which is finite over $X$, and $Z$ by its inverse image in $X'$, we may assume
$X$ normal.

\medskip \textit{Step 2}. \textit{Use of \ref{10-de Jong-Gabber-Vidal1}}.
Choose a finite Galois extension $k_0$ of $k$ such that the irreducible
components of $X_0 = X \times_S S_0$ ($S_0 = \Spec \, k_0$) are geometrically
irreducible. Let $G = \mathrm{Gal}(k_0/k)$ and $H \subset G$ the
decomposition subgroup of a component $Y_0$ of $X_0$. We apply \ref{10-de
Jong-Gabber-Vidal1} to ($Y_0/S_0$, $Z_0 \cap Y_0$), where $Z_0 = S_0 \times_S
Z$.  We find a surjection $H_1 \ra H$ and an $H_1$-equivariant diagram of
type \ref{10-Vidal-diagram1} :
\begin{equation}\label{10-Vidal-diagram2}
\xymatrix{Y_0 \ar[d] & Y_1 \ar[d] \ar[l]_{b'} & Y_2 \ar[l]_{a'} \ar[dl]
\\ S_0 & S_1 \ar[l] & {}},
\end{equation}
satisfying conditions (i), (ii),
(iii), (iv) with $S$ replaced by $S_0$, and $X_2 \ra X_1 \ra X$ by $Y_2 \ra
Y_1 \ra Y_0$. As $G$ transitively permutes the components of $X_0$, $X_0$ is,
as a $G$-scheme over $S_0$, the contracted product $$ X_0 = Y_0 \times^H G,
$$ i. e. the quotient of $Y_0 \times G$ by $H$ acting on $Y_0$ on the right
and on $G$ on the left (cf. proof of \refext{VIII}{8-rigidlem2}), and similarly
$Z = Z_0\times^H G$. Choose an extension of the diagram $\xymatrix{H_1 \ar[r]^{u}
& H \ar[r]^{i} & G}$ into a commutative diagram of finite groups $$
\xymatrix{H_1 \ar[d]^{u} \ar[r]^{i_1} & G_1 \ar[d]^v \\ H \ar[r]^i & G} $$
with $i_1$ injective and $v$ surjective (for example, take $i_1$ to be the
graph of $iu$ and $v$ the projection). Define $$ X_1 := Y_1 \times^{H_1} G_1,
\, \, \, X_2 := Y_2 \times^{H_1} G_1.  $$ Then (\ref{10-Vidal-diagram2})
extends to a $G_1$-equivariant diagram of type \ref{10-Vidal-diagram1}
\begin{equation}\label{10-Vidal-diagram3} \xymatrix{X_0 \ar[d] & X_1 \ar[d]
\ar[l]_b & X_2 \ar[l]_a \ar[dl] \\ S_0 & S_1 \ar[l] & {}},
\end{equation}
satisfying again (i), (ii), (iii), (iv). In particular, the composition $h
: X_2 \ra X_1 \ra X_0 \ra X$ is an alteration above the composition $S_1
\ra S_0 \ra S$, $X_2$ is projective and smooth over $S_1$, and $h^{-1}(Z)$
is the support of an snc divisor. However, as regard to \ref{10-de Jong1},
the diagram $$ \xymatrix{X \ar[d] & X_2 \ar[l]_h \ar[d] \\ S & S_1 \ar[l]}
$$ deduced from (\ref{10-Vidal-diagram3}) has two defects :

(a) the extension $S_1 \ra S$ is not necessarily of degree prime to $\ell$,

(b) the alteration $h$ is not necessarily an $\ell'$-alteration.

We will first repair (a) and (b) at the cost of temporarily losing the
smoothness of $X_2/S_1$ and the snc property of $h^{-1}(Z)$. By (i),
$S_1/\mathrm{Ker}(G_1 \ra G) \ra S_0$ is a radicial extension, hence
$S_1/G_1 \ra S = S_0/G$ is a radicial extension, too. Similarly, by (iii),
$X_2/G_1 \ra X$ is an alteration over $S_1/G_1 \ra S$,  which is a universal
homeomorphism over a dense open subset. Now let $L$ be an $\ell$-Sylow subgroup
of $G_1$. Then $S_1/L \ra S_1/G_1$ is of degree prime to $\ell$, and $X_2/L \ra
X_2/G_1$ is a finite surjective morphism of generic degree prime to $\ell$. Let
$S' := \Spec \, k' = S_1/L$, $X' = S' \times_S X$. We get a commutative diagram
with cartesian square
\begin{equation}\label{10-Vidal-diagram4} \xymatrix{X
\ar[d] & X' \ar[d] \ar[l] & X_2/L \ar[l] \ar[dl] & X_2 \ar[l] \ar[dl]\\ S &
S' \ar[l] & S_1\ar[l]},
\end{equation}
where $S' \ra S$ is an extension of
degree prime to $\ell$, $S_1 \ra S'$ a Galois extension of group $L$, $$ h_2
: X_2/L \ra X $$ an $\ell'$-alteration, $X_2/S_1$ is projective and smooth,
and if now $h$ denotes the composition $X_2 \ra X$, $Z_1 := h^{-1}(Z)$ is
an snc divisor in $X_2$. If $X_2/L$ was smooth over $S'$ and $Z_1/L$ an snc
divisor in $X_2/L$, we would be finished. However, this is not the case in
general. We will use Gabber's theorem \ref{10-modif-logsmooth}  to fix this.

\medskip \textit{Step 3.} \textit{Use of \ref{10-modif-logsmooth}.} Let $Y$
be a connected component of $X_2$, $(Z_1)_Y = h^{-1}(Z) \cap Y$, $D$ the
stabilizer of $Y$ in $L$, $I \subset D$ the inertia group at the generic
point of $Y$. Then $D$ acts on $Y$ through $K := D/I$, and this action
is generically free. As $Y$ is smooth over $S_1$ and $(Z_1)_Y$ is an snc
divisor in $Y$, $(Y,(Z_1)_Y)$ makes a log regular pair, log smooth over $S_1$,
hence over $S' = S_1/L$ (equipped with the trivial log structure). We have a
$K$-equivariant commutative diagram
\begin{equation}\label{10-Gabber-diagram1}
\xymatrix{(Y/K,(Z_1)_Y/K) \ar[d] & (Y,(Z_1)_Y) \ar[l] \ar[dl]_f \\ S'},
\end{equation} where $K$ acts trivially on $S'$, and $f$ is projective,
smooth, and log smooth ($S'$ having the trivial log structure). We now apply
\ref{10-modif-logsmooth} to ($f : (Y,(Z_1)_Y) \ra S' = (S',\emptyset),
K$), which satisfies conditions (i) - (iii) of (\textit{loc. cit.}). We
get a $D$-equivariant projective modification $g : Y_1 \ra Y$ ($D$ acting
through $K$) and a $D$-strict snc divisor $E_1$ on $Y_1$, containing
$Z_1 := g^{-1}((Z_1)_Y)$ as a subdivisor, such that the action of $D$
on $(Y_1,E_1)$ is very tame, and $(Y_1,E_1)$ and $(Y_1/D,E_1/D)$ are log
smooth over $S'$. Pulling back $g$ to the orbit $Y \times^D L$ of $Y$
under $L$, i. e. replacing $g$ by $g \times^D L$, and working separately
over each orbit, we eventually get an $L$-equivariant commutative square
\begin{equation}\label{10-Gabber-diagram2} \xymatrix{(Y_2/L,E_2/L) \ar[d]_v &
(Y_2,E_2) \ar[l] \ar[d]_u\\ (X_2/L,Z_1/L) & (X_2,Z_1) \ar[l]},
\end{equation}
where $u$, $v$ are projective modifications (and $Z_1 = h^{-1}(Z)$, $Z_1/L =
h_2^{-1}(Z)$ as above), with the property that the pair $(Y_2/L,E_2/L)$ is
an fs log scheme log smooth over $S'$ ($= S_1/L$), and $v^{-1}(h_2^{-1}(Z))
\subset E_2/L)$. Let $w : (\tilde X, \tilde E) \ra (Y_2/L,E_2/L)$ be a
projective, log étale modification such that $\tilde X$ is regular, and
$\tilde E = w^{-1}(E_2/L)$ is an snc divisor in $\tilde X$. For example,
one can take for $w$ the saturated monoidal desingularization $\tilde
{\mathscr{F}}^{\mathrm{log}}$ of (\refext{VIII}{8-deslogregth}). We then apply
\ref{10-local-structure} to the log smooth morphism $\tilde X \ra S'$. By a
special case of the (\ref{10-chart-diagram}), with $Q = \{1\}$, $G = \{1\}$,
as $P^*$ is torsion free, $\tilde X$ is not only regular, but smooth over $S'$,
and $\tilde E$ a relative snc on $\tilde X$. Let $$ \tilde h : \tilde X \ra X
$$ be the composition $$ \xymatrix{\tilde X \ar[r]^w & Y_2/L \ar[r]^v & X_2/L
\ar[r]^{h_2} & X}.  $$ This is a projective $\ell'$-alteration, and it fits
in the commutative diagram $$ \xymatrix{X \ar[d] & \tilde X \ar[l]_{\tilde
h} \ar[d] \\ S & S' \ar[l]}, $$ where $S'$ is an extension of $S$ of degree
prime to $\ell$, $\tilde X$ is projective and smooth over $S'$, and $\tilde
h^{-1}(Z)_{\mathrm{red}}$ is a sub-divisor of the $S'$-relative snc divisor
$\tilde E$, hence a relative snc divisor as well. This finishes the proof
of \ref{10-de Jong1}.

Recall now the theorem stated in (Introduction 0.3 (2)) :
\begin{theorem}[Gabber]\label{10-de Jong2} Let $S$ be a separated, integral,
noetherian, excellent, regular scheme of dimension 1, with generic point
$\eta$, $X$ a scheme separated, flat and of finite type over $S$, $\ell$
a prime number invertible on $S$, $Z \subset X$ a nowhere dense closed
subset. Then there exists a finite extension $\eta'$ of $\eta$ of degree prime
to $\ell$ and a projective $\ell$'-alteration $h : \tilde X \ra X$ above $S'
\ra S$, where $S'$ is the normalization of $S$ in $\eta'$, with $\tilde X$
regular and quasi-projective over $S'$, a strict normal crossings divisor
$\tilde T$ on $\tilde X$, and a finite closed subset $\Sigma$ of $S'$ such
that :

(i) outside $\Sigma$, $\tilde X \ra S'$ is smooth and $\tilde T \ra S'$
is a relative divisor with normal crossings ;

(ii) étale locally around each geometric point $x$ of $\tilde X_{s'}$,
where $s' = \Spec \, k'$ belongs to $\Sigma$, the pair $(\tilde X, \tilde T)$
is isomorphic to the pair consisting of
$$ X'=S'[u_1^{\pm 1}, \cdots, u_s^{\pm 1}, t_1, \cdots, t_n]/
(u_1^{b_1} \cdots u_s^{b_s}t_1^{a_1} \cdots t_r^{a_r}- \pi), $$
$$ T'=V(t_{r+1}\cdots t_m)\subset X'$$
around the point $(u_i = 1), (t_j = 0)$, with $1
\le r \le m \le n$, for positive integers $a_1, \cdots, a_r, b_1, \cdots, b_s$
satisfying $\mathrm{gcd}(p,a_1, \cdots, a_r, b_1, \cdots, b_s) = 1$, $p$ the
characteristic exponent of $k'$, $\pi$ a local uniformizing parameter at $s'$ ;

(iii) $\tilde h^{-1}(Z)_{\mathrm{red}}$ is a sub-divisor of $\cup_{s' \in
\Sigma}(\tilde X_{s'})_{\mathrm{red}} \cup \tilde T$.
\end{theorem}

The proof follows the same lines as that of \ref{10-de Jong1}. We
need again a Gabber-Vidal variant of de Jong's alteration
theorems (cf. \cite[3.8]{Independance@Zheng}). This is essentially
\cite[4.4.1]{BS@Vidal}), except for the additional data of $Z \subset X$
and the removal of the hypothesis that $S$ is a strictly local trait :

\begin{lemma}\label{10-de Jong-Gabber-Vidal2} Let $X$ be a normal,
proper scheme over $S$, whose generic fiber is geometrically reduced and
irreducible, $Z \subset X$ a nowhere dense closed subset. We assume that
a finite group $H$ acts on $X \ra S$, faithfully on $X$, and that $Z$ is
$H$-stable. Then there exists a finite group $H_1$, a surjective homomorphism
$H_1 \ra H$, and an $H_1$-equivariant diagram with a cartesian square
\begin{equation}\label{10-Vidal-diagram5} \xymatrix{X \ar[d] & X_1 \ar[d]
\ar[l]_b & X_2 \ar[l]_a \ar[dl] \\ S & S_1 \ar[l] & {}},
\end{equation}
satisfying the following properties :

(i) $S_1 \ra S$ is the normalization of $S$ in a finite extension $\eta_1$
of $\eta$ such that $\eta_1/(\mathrm{Ker}(H_1 \ra H) \ra \eta$ is a radicial
extension (where $\eta$ is the generic point of $S$) ;

(ii) $X_2$ is projective and strictly semistable over $S_1$ (i. e. is
strictly semistable over the localizations of $S_1$ at closed points
\cite[2.16]{Alterations@de_Jong}) ;

(iii) $a : X_2 \ra X_1$ is  projective and surjective, maximally dominating and
generically finite and flat, and  there exists an $H_1$-admissible dense open
subset $W \subset X_2{}_{\eta_1}$ over a dense open subset $U$ of $X_{\eta}$,
such that if $U_1 = \eta_1 \times_{\eta} U$ and $K = \mathrm{Ker}(H_1 \ra
\mathrm{Aut}(U_1))$, $W \ra W/K$ is a Galois étale cover of group $K$
and the morphism $W/K \ra U_1$ induced by $a$ is a universal homeomorphism ;

(iv) $(ba)^{-1}(Z)$ is the support of a strict normal crossings
divisor in $X_2$, and $(X_2,(ba)^{-1}(Z))$ is a strict semistable pair
over $S_1$ (i. e. over the localizations of $S_1$ at closed points
\cite[6.3]{Alterations@de_Jong}).
\end{lemma}

\medskip Note that (ii) and (iv) imply that there exists a finite closed subset
$\Sigma$ of $S_1$ such that, outside $\Sigma$, the pair $(X_2,(ba)^{-1}(Z))$
consists of a smooth morphism and a relative strict normal crossings divisor.

\medskip \textit{Proof.} Up to minute modifications the proof is the
same as that of \ref{10-de Jong-Gabber-Vidal1}. We may assume the generic
fiber $X_{\eta}$ of dimension $d \ge 1$. We apply \cite[4.4.3]{BS@Vidal}
to $X/S$, $Z$, and $G = H$. We get the data of (\textit{loc. cit.}), namely
an equivariant finite surjective morphism $(S_1, H_1) \ra (S,H)$, with $S_1$
regular (hence equal to the normalization of $S$ in the extension $\eta_1$
of the generic point $\eta$ of $S$) such that $\eta_1/\mathrm{Ker}(H_1 \ra H)
\ra \eta$ is radicial, an $H_1$-equivariant pluri-nodal fibration $(Y_d \ra
\cdots \ra Y_1 \ra S_1, \{ \sigma_{ij} \}, Z_0)$, and an $H_1$-equivariant
alteration $a_1 : Y_d \ra X$ over $S$, satisfying the conditions (i),
(ii), (iii) of (\textit{loc. cit.}) (in particular $a_1^{-1}(Z) \subset
Z_d$). Then, as in the proof of \cite[4.4.1]{BS@Vidal}, successively applying
\cite[4.4.4]{BS@Vidal} to each nodal curve $f_i : Y_i \ra Y_{i-1}$, one
can replace $Y_i$ by an $H_1$-equivariant projective modification $Y'_i$
of it such that $Y'_i$ is regular, and the inverse image of $Z_i := \cup_j
\sigma_{ij}(Y_{i-1}) \cup f_i^{-1}(Z_{i-1})$ in $Y'_i$ is an $H_1$-equivariant
strict snc divisor. Then, by the remark following \cite[4.4.4]{BS@Vidal} $X_2
:= Y'_d$ is strict semistable over $S_1$ and $(X_2,Z_d)$ is a strict semistable
pair over $S_1$. In particular, $(ba)^{-1}(Z)_{\mathrm{red}}$ is a subdivisor
of $Z_d$, hence an snc divisor, and $(X_2,(ba)^{-1}(Z))_{\mathrm{red}}$ is
a strict semistable pair over $S_1$. The open subsets $U$ and $V$ of (iii)
are constructed as at the end of the proof of \cite[4.4.1]{BS@Vidal}.

\subsection{}\label{10-proof of de Jong2} \textit{Proof of \ref{10-de
Jong2}}. It is similar to that of \ref{10-de Jong1}. There are again three
steps. We will indicate which modifications should be made.

\medskip \textit{Step 1.} \textit{Preliminary reductions}. Up to replacing $X$
by the disjoint union of the schematic closures of the reduced components of
its generic fiber, and working separately with each of them, we may assume $X$
integral (and $X_{\eta} \ne \emptyset$). Applying Nagata's compactification
theorem, we may further assume $X$ proper and integral.  Base changing by
the normalization of $S$ in a finite radicial extension of $\eta$ and taking
the reduced scheme, we reduce to the case where, in addition, $X_{\eta}$
is irreducible and geometrically reduced. Then we blow up $Z$ in $X$ and
normalize as in the previous step 1. Here we used the excellency of $S$
to guarantee the finiteness of the normalizations.

\medskip \textit{Step 2.} \textit{Use of \ref{10-de Jong-Gabber-Vidal2}}. Let
$S_0$  be the normalization of $S$ in a finite Galois extension $\eta_0$ of
$\eta$ such that the irreducible components of the generic fiber of $X_0 =
X \times_S S_0$ ($S_0 = \Spec \, k_0$) are geometrically irreducible. Let $G
= \mathrm{Gal}(\eta_0/\eta)$ and $H \subset G$ the decomposition subgroup
of a component $Y_0$ of $X_0$. We apply \ref{10-de Jong-Gabber-Vidal2}
to ($Y_0/S_0$, $Z_0 \cap Y_0$), where $Z_0 = S_0 \times_S Z$.  We find
a surjection $H_1 \ra H$ and an $H_1$-equivariant diagram of type
(\ref{10-Vidal-diagram5}) satisfying conditions (i), (ii), (iii), (iv)
with $S$ replaced by $S_0$, and $X_2 \ra X_1 \ra X$ by $Y_2 \ra Y_1 \ra
Y_0$. We then, as above,  extend $H_1 \ra H$ to a surjection $G_1 \ra G$
and obtain a $G_1$-equivariant diagram of type (\ref{10-Vidal-diagram5})
\begin{equation}\label{10-Vidal-diagram6} \xymatrix{X_0 \ar[d] & X_1 \ar[d]
\ar[l]_b & X_2 \ar[l]_a \ar[dl] \\ S_0 & S_1 \ar[l] & {}},
\end{equation}
satisfying again (i), (ii), (iii), (iv). In particular, the composition $h :
X_2 \ra X_1 \ra X_0 \ra X$ is an alteration above the composition $S_1 \ra S_0
\ra S$, $X_2$ is projective over $S_1$, with strict semistable reduction, and
$h^{-1}(Z)$ is the support of an snc divisor forming a strict semistable pair
with $X_2/S_1$. It follows from (i) that $S_1/G_1 \ra S = S_0/G$ is generically
radicial, and by (iii) that $X_2/G_1 \ra X$ is an alteration over $S_1/G_1 \ra
S$,  which is a universal homeomorphism over a dense open subset. As above,
choose an $\ell$-Sylow subroup $L$ of $G_1$. Then $S_1/L$ is regular, $S_1/L
\ra S_1/G_1$ is finite surjective of generic degree prime to $\ell$, and
$X_2/L \ra X_2/G_1$ is a finite surjective morphism of generic degree prime
to $\ell$. Putting $S' = S_1/L$, $X' = S' \times_S X$, we get a commutative
diagram with cartesian square
\begin{equation}\label{10-Vidal-diagram7}
\xymatrix{X \ar[d] & X' \ar[d] \ar[l] & X_2/L \ar[l] \ar[dl] & X_2 \ar[l]
\ar[dl]\\ S & S' \ar[l] & S_1\ar[l]},
\end{equation}
where $S'$ is regular,
$S' \ra S$ is finite surjective of generic degree prime to $\ell$, $S_1 \ra
S'$ generically étale of degree the order of $L$, $$ h_2 : X_2/L \ra X $$
an $\ell'$-alteration, $X_2/S_1$ is projective and strictly semistable, and
if $h$ denotes the composition $X_2 \ra X$, $Z_1 := h^{-1}(Z)_{\mathrm{red}}$
is an snc divisor in $X_2$, forming a strictly semistable pair with $X_2/S_1$.

\medskip \textit{Step 3}. \textit{Use of \ref{10-modif-logsmooth}}. Defining
$Y$, $(Z_1)_Y$, $I \subset D$, $K = D/I$ as in the former step 3, $K$ acts
generically freely on $Y$. As the pair $(Y,(Z_1)_Y)$ is strictly semistable
over $S_1$, there exists a finite closed subset $\Sigma_1$ of $S_1$ such that
$(Y, Y_{\Sigma_1} \cup (Z_1)_Y)$ forms a log regular pair, log smooth over
$S_1$ equipped with the log structure defined by $\Sigma_1$. As $S_1 \ra S'
= S_1/L$ is Kummer étale, $(Y,(Y_{\Sigma'})_{\mathrm{red}} \cup (Z_1)_Y)$
(where $\Sigma'$ is the inverse image of $\Sigma_1$) is also log smooth
over $S'$ (equipped with the log structure given by $\Sigma'$), and we get a
$K$-equivariant commutative diagram (\ref{10-Gabber-diagram1}), with trivial
action of $K$ on $S'$ and $f$ projective and log smooth over $S'$. We then
apply \ref{10-modif-logsmooth} to $f : (Y,(Y_{\Sigma'})_{\mathrm{red}} \cup
(Z_1)_Y) \ra S'$, and the proof runs as above. We get a $D$-equivariant
projective modification $g : Y_1 \ra Y$ ($D$ acting through $K$) and a
$D$-strict snc divisor $E_1$ on $Y_1$, containing $(g^{-1}((Z_1)_Y) \cup
(Y_1)_{\Sigma'})_{\mathrm{red}}$ as a subdivisor, such that the action
of $D$ on $(Y_1,E_1)$ is very tame, and $(Y_1,E_1)$ and $(Y_1/D,E_1/D)$
are log smooth over $S'$. After extending from $D$ to $L$ we get an
$L$-equivariant commutative square of type (\ref{10-Gabber-diagram2}), with
$(Y_2/L,E_2/L)$ log smooth over $S'$ ($= S_1/L$), and $(v^{-1}(h_2^{-1}(Z))
\cup (Y_2/L)_{\Sigma'})_{\mathrm{red}} \subset E_2/L$. As above, we take
a projective, log étale modification such that $\tilde X$ is regular,
and $\tilde E = w^{-1}(E_2/L)$ is an snc divisor in $\tilde X$.

We now apply \ref{10-local-structure} to the log smooth morphism
$(\tilde X, \tilde E) \ra S'$. It's enough to work étale locally on
$\tilde X$ around some geometric point $x$ of $\tilde X_{s'}$, with $s'
\in \Sigma'$. We replace $S'$ by its strict localization at the image
of $x$, and consider (\ref{10-chart-diagram}), with $Q = \mathbf{N}$,
$G = \{1\}$, $\Lambda = \mathbf{Z}_{(p)}$ if $p > 1$ and $\mathbf{Q}$
otherwise, and the chart $a : \mathbf{N} \ra M_{S'}, 1 \mapsto \pi$,
$\pi$ a uniformizing parameter of $S'$. In (\ref{10-P1}) we have $P^*=
\mathbf{Z}^{\mu}$, $P_1 = \mathbf{N}^{\nu}$, for nonnegative integers
$\mu$, $\nu$, hence $$ P = \mathbf{Z}^{\nu} \oplus \mathbf{N}^{\mu}.
$$ Let $((b_i),(a_i)$ be the image of $1 \in \mathbf{N}$ in $P$ in the
above decomposition, and let $(a_1,\cdots,a_r)$, $(b_1,\cdots,b_s)$ be the
sets of those $a_i$'s and $b_i$'s which are $\ne 0$. We may assume $b_i
>0$ if $b_i \ne 0$. As the torsion part of $\mathrm{Coker}(\mathbf{Z}
\ra P^{gp})$ is annihilated by an integer invertible on $\tilde X$, we
have $\mathrm{gcd}(p,a_1,\cdots,a_r,b_1,\cdots,b_s) = 1$, where $p$ is the
characteristic exponent of $k$. We have $P = \mathbf{Z}^s \oplus \mathbf{N}^r
\oplus \mathbf{Z}^{\nu-s} \oplus \mathbf{N}^{\mu - r}$. Choosing a basis of
$V$, we get that étale locally around $x$, $\tilde X$ is given by a cartesian
square $$ \xymatrix{\tilde X \ar[d] \ar[r] & \Spec \, \Lambda[u_1^{\pm 1},
\cdots, u_s^{\pm 1}, t_1, \cdots, t_n] \ar[d]\\ S' \ar[r] & \Spec \, \Lambda[z]}
$$ with $x$ going to the point $(u_i = 1), (t_j = 0)$, and $z \mapsto \pi$,
$z \mapsto u_1^{b_1} \cdots u_s^{b_s}t_1^{a_1}\cdots t_r^{a_r}$, in other
words, $$ \tilde X = S'[u_1^{\pm 1}, \cdots, u_s^{\pm 1}, t_1, \cdots,
t_n]/(u_1^{b_1} \cdots u_s^{b_s}t_1^{a_1} \cdots t_r^{a_r} - \pi), $$ Finally,
$\tilde E$ is the union of the special fiber $\tilde X_{s'}$ and a horizontal
divisor $\tilde T$, étale locally given by the equation $t_{r+1}\cdots
t_m = 0$, for an integer $m$ such that $1 \le r \le m \le n$. As $\tilde
h^{-1}(Z)_{\mathrm{red}}$ is a sub-divisor of $(\tilde X_{s'})_{\mathrm{red}}
\cup \tilde T$, this finishes the proof of \ref{10-de Jong2}.

\section{Resolvability, log smoothness, and weak semistable reduction}\label{10-complsec}

\subsection{Elimination of separatedness assumptions}\label{10-sepelimsec}
The main aim of \S\ref{10-sepelimsec} is to weaken the separatedness assumptions
in Theorems \refext{VIII}{8-TH. MODIFICATION} and \ref{10-modif-logsmooth}.
We start with an example showing how things can go wrong without such an assumption.
Recall, see \refext{VI}{6-stratification}, that if a finite group $G$ acts on
a scheme $X$ then the fixed point subscheme $X^G$ is the intersection of graphs
of the translations $g{\colon}X\to X$.  In particular, $X^G$ is closed whenever $X$
is separated.

\begin{example2}\label{10-nonsepex}
Let $G=\{1,\sigma\}$ be a cyclic group of order two. Consider an affine plane
$\AA^2_k=\Spec(k[x,y])$ over a field $k$ of characteristic different from $2$
and let $X$ be the affine plane with doubled origin, i.e., $X$ is glued from two
copies of $\AA^2_k$ along $U=\AA^2_k\setminus\{0\}$. We provide $U$ with the
action of $G$ by the rules $\sigma x=x$ and $\sigma y=-y$, and extend this action
to $X$ so that $\sigma$ intertwines the origins. Then $X^G$ is the closed subscheme
of $U$ given by $y=0$. In particular, it is not closed in $X$, and the inertia
stratification (used in the proof of Theorem \refext{VIII}{8-TH. MODIFICATION})
does not make sense. Note that the closure $Z$ of $X^G$ is an affine line with
doubled origin on which $G$ acts by permuting the origins. For any $G$-equivariant
modification $Y\to X$ the proper transform of $Z$ is a modification of $Z$ and hence
is isomorphic to $Z$. Thus, $Y^G$ contains the generic point of $Z$ but not its origins,
and we obtain that $Y^G$ is not closed. Therefore, the action on $Y$ is not very tame
for any choice of a log structure, and we obtain that the assertion of Theorem
\refext{VIII}{8-TH. MODIFICATION} fails for such $(G,X,Z=\emptyset)$.
\end{example2}

\subsubsection{Inertia specializing actions}
Example \ref{10-nonsepex} motivates the following definition: an action of
a finite group $G$ on a scheme $X$ is {\em inertia specializing} if for any
point $x\in X$ with a specialization $y\in X$ one has that $G_x\subseteq G_y$.
Since a subscheme of $X$ is closed if and only if it is closed under specializations,
we obtain the following result.

\begin{lemma2}
An action of $G$ on $X$ is inertia specializing if and only if for each subgroup
$H\subseteq G$ the subscheme $X^H$ is closed.
\end{lemma2}

\begin{remark2}\label{10-inspecrem}
(i) A large class of examples of inertia specializing actions can be described as follows.
The following conditions are equivalent and imply that the action is inertia specializing:
(a) any $G$ orbit is contained in an open separated subscheme of $X$, (b) $X$ admits
a covering by $G$-equivariant separated open subschemes $X_i$. In particular, any
admissible action is inertia specializing.

(ii) If $(G,X,Z)$ is as in Theorem \refext{VIII}{8-TH. MODIFICATION}, but instead of
separatedness of $X$ one only assumes that it possesses a covering by $G$-equivariant
separated open subschemes $X_i$, then the assertion of the theorem still holds true.
Indeed, the theorem applies to the $G$-equivariant log schemes $(X_i,Z_i=Z|_{X_i})$,
and by Theorem \refext{VIII}{8-addendum} the modifications $f_{(G,X_i,Z_i)}$ agree
on the intersections and hence glue to a required modification $f_{(G,X,Z)}$ of $X$.
\end{remark2}

A quick analysis of the proof of \refext{VIII}{8-TH. MODIFICATION} is required
to obtain the following stronger result.

\begin{theorem2}\label{10-elimth}
(i) Theorem \refext{VIII}{8-TH. MODIFICATION} and its complement \refext{VIII}{8-addendum}
hold true if the assumption that $X$ is separated is weakened to the assumption that
the action of $G$ on $X$ is inertia specializing.

(ii) Theorem \ref{10-modif-logsmooth} holds true if the assumptions that $X$ and $S$ are
separated are replaced with the single assumption that the action of $G$ on $X$ is inertia specializing.
\end{theorem2}
\begin{proof}
The proof of \refext{VIII}{8-TH. MODIFICATION} in the general (i.e. not necessarily qe)
case runs as follows. First, for any point $x\in X$ and localizations $X_x$ and $Z_x$
one constructs $f_x=f_{(G_x,X_x,Z_x)}$. This stage does not use any separatedness assumption.
Second, for any specialization $y$ of $x$ one shows that after eliminating empty normalized
blow ups, the restriction  $f_y|_{X_x}$ coincides with $f_x$. For this stage to work we only
need that $G_x\subseteq G_y$, i.e. the assumption that the action is inertia specializing is
precisely what one uses. The third stage is to show that the family $f_y$ forms an atomic
normalized blow up tower and hence descends to a normalized blow up tower $f_{(G,X,Z)}$.
It uses the log-inertia stratification in order to control the associated points of the
normalized blow up towers $f_x$. Since the inertia stratification is well defined for
inertia specializing actions, the log-inertia stratification is well defined too, and
this stage works fine. Finally, the last stage is to uniformly bound the normalization
threshold of $f_{(G,X,Z)}$ in terms of the combinatorial datum. The argument is local
and the combinatorial datum involves only the log-inertia stratification and the size
of $G$. Hence this stage also extends to our case without any changes.

The proof of Theorem \ref{10-modif-logsmooth} runs as follows. One considers the modification
$f_{(G,X,Z)}$ from \refext{VIII}{8-TH. MODIFICATION} and checks that it satisfies the additional
properties asserted by \ref{10-modif-logsmooth}. This check is local on $X$ and hence applies to
non-separated schemes as well. Since  by part (i) of \ref{10-elimth}, $f_{(G,X,Z)}$ is well defined
whenever $G$ acts inertia specializing on $X$, we obtain (ii).
\end{proof}

\subsubsection{Pseudo-projective morphisms and non-separated Chow's lemma}\label{10-chowsec}
We conclude \S\ref{10-sepelimsec} with recalling a non-separated version of Chow's lemma
due to Artin-Raynaud-Gruson. It will be needed to avoid unnecessary separatedness
assumptions in the future. We prefer to use the following non-standard terminology:
a finite type morphism $f{\colon}X\to S$ is {\em pseudo-projective} if it can be factored
into a composition of a {\em local isomorphism} $X\to\oX$ (i.e. $X$ admits an open covering
$X=\cup_i X_i$ such that the morphisms $X_i\to\oX$ are open immersions) and a projective morphism $\oX\to S$.

\begin{remark2}\label{10-chowrem}
(i) We introduce pseudo-projective morphisms mainly for terminological convenience.
Although pseudo-projectivity is preserved by base changes, it can be lost under
compositions. Moreover, even if $X$ is pseudo-projective over a field $k$, its
blow up $X'$ does not have to be pseudo-projective over $k$ (thus giving an example
of a projective morphism $f{\colon}X'\to X$ and a pseudo-projective one $X\to\Spec(k)$
so that the composition is not pseudo-projective). Indeed, let $X$ be an affine plane
with a doubled origin $\{o_1,o_2\}$, and let $X'$ be obtained by blowing up $o_1$. By
$\eta$ we denote the generic point of $C_1=f^{-1}(o_1)$. The ring $\Spec(\calO_{X',\eta})$
is a DVR and its spectrum has two different $k$-morphisms to $X'$: one takes the closed point
to $\eta$ and another one takes it to $o_2$. It then follows from the valuative criterion of
separatedness that any $k$-morphism $g{\colon}X'\to Y$ with a separated target takes $o_2$
and $\eta$ to the same point of $Y$. In particular, such $g$ cannot be a local isomorphism,
and hence $X'$ is not pseudo-projective over $k$.

(ii) Note that a morphism $f$ is separated (resp. proper) and pseudo-projective if and only
if it is quasi-projective (resp. projective). So, the following result extends the classical
Chow's lemma to non-separated morphisms.
\end{remark2}

\begin{proposition2}\label{10-chowlem}
Let $f{\colon}X\to S$ be a finite type morphism of quasi-compact and quasi-separated schemes,
and assume that $X$ has finitely many maximal points. Then there exists a projective
modification $g{\colon}X'\to X$ (even a blow up along a finitely generated ideal with
a nowhere dense support) such that the morphism $X'\to S$ is pseudo-projective.
\end{proposition2}
\begin{proof}
As a simple corollary of the flattening theorem, it is proved in \cite[5.7.13]{Platification@Raynaud-Gruson}
that there exists a modification $X'\to X$ such that $X'\to S$ factors as a composition of an
étale morphism $X'\to\oX$ that induces an isomorphism of dense open subschemes and a projective
morphism $\oX\to S$. (In loc.cit. one works with algebraic spaces and assumes that $f$ is locally
separated, but the latter is automatic for any morphism of schemes.) Our claim now follows from
the following lemma (which fails for locally separated morphisms between algebraic spaces).
\end{proof}

\begin{lemma2}
Assume that $\phi{\colon}Y\to Z$ is an étale morphism of schemes that restricts to an open
embedding on a dense open subscheme $Y_0\into Y$. Then $\phi$ is a local isomorphism.
\end{lemma2}
\begin{proof}
Let us prove that if, in addition, $\phi$ is separated then $\phi$ is an open immersion.
Since $Y$ possesses an open covering by separated subschemes, this will imply the lemma.
The diagonal $\Delta_\phi{\colon}Y\to Y\times_ZY$ is an open immersion, and by the
separatedness of $\phi$, it is also a closed immersion. Thus, $Y$ is a connected
component of $Y\times_ZY$, and since both $Y$ and $Y\times_ZY$ have dense open
subschemes that map isomorphically onto $Y_0$, $\Delta_\phi$ is an isomorphism.
This implies that $\phi$ is a monomorphism, but any étale monomorphism is an
open immersion by \cite[17.9.1]{EGAIV-4}.
\end{proof}

\subsection{Semistable curves and log smoothness}

\subsubsection{Log structure associated to a closed subset}
Let $S$ be a reduced scheme. Any closed nowhere dense subset $W\subset S$ induces
a log structure $j_*\calO^*_U\cap\calO_S\into\calO_S$, where $j{\colon}U\into S$
is the embedding of the complement of $W$. The associated log scheme will be denoted
$(S,W)$. By \refext{VI}{6-1.4}, any log regular log scheme is of the form $(S,W)$,
where $W$ is the non-triviality locus of the log structure.

\subsubsection{Semistable relative curves}\label{10-relcurvesec}
Following the terminology of \cite{Temkin@stable}, by a semistable multipointed
relative curve over a scheme $S$ we mean a pair $(C,D)$, where $C$ is a flat
$S$-scheme of relative dimension one and with fibers having only ordinary nodes
as singularities, and $D\into C$ is a closed subscheme which is étale over $S$
and disjoint from the singular locus of $C\to S$. We do not assume $C$ to be neither
proper nor even separated over $S$.

\begin{proposition2}\label{10-logsmprop}
Assume that $(S,W)$ is a log regular log scheme and $(C,D)$ is a semistable multipointed
relative $S$-curve such that the morphism $f{\colon}C\to S$ is smooth over $S\setminus W$.
Then the morphism of log schemes $(C,D\cup f^{-1}(W))\to(S,W)$ is log smooth.
\end{proposition2}
\begin{proof}
If $C$ is proper over $S$ then this claim is nothing else than \refext{VI}{6-nodal implique
log lisse}. Moreover, the proof of \refext{VI}{6-nodal implique log lisse}. is local on $C$,
hence it applies to our case as well.
\end{proof}

\subsection{$l'$-resolvability}\label{10-lresolvesec}

\subsubsection{Alterations}
Assume that $S'$ and $S$ are reduced schemes with finitely many maximal points and let
$\eta'\subseteq S'$ and $\eta\subseteq S$ denote the subschemes of maximal points.
Let $f{\colon}S'\to S$ be an {\em alteration}, i.e. a proper, surjective, generically finite,
and maximally dominating morphism. Recall that $f$ is an {\em $l'$-alteration} if
$([k(x):k(f(x))],l)=1$ for any $x\in\eta'$. We say that $f$ is {\em separable} if
$k(\eta')$ is a separable $k(\eta)$-algebra (i.e. $\eta'\to\eta$ has geometrically
reduced fibers). If, in addition, $S'$ and $S$ are provided with an action of
a finite group $G$ such that $f$ is $G$-equivariant, the action on $S$ is trivial,
the action on $\eta'$ is free, and $\eta'/G\iso\eta$, then we say that $f$ is a
{\em separable Galois alteration} of group $G$, or just {\em separable $G$-alteration}.

\begin{remark2}
We add the word "separable" to distinguish our definition from Galois alterations
in the sense of de Jong (see \cite{Alterations@de_Jong}) or Gabber-Vidal
(see \cite[p. 370]{Brauer@Vidal}). In the latter cases, one allows alterations
that factor as $S'\to S''\to S$, where $S'\to S''$ is a separable Galois alteration
and $S''\to S$ is purely inseparable.
\end{remark2}

\subsubsection{Universal $l'$-resolvability}\label{10-resolvsec}
Let $X$ be a noetherian scheme and let $l$ be a prime invertible on $X$. Assume that
for any alteration $Y\to X$ and nowhere dense closed subset $Z\subset Y$ there exists
a surjective projective morphism $f{\colon}Y'\to Y$ such that $Y'$ is regular and
$Z'=f^{-1}(Z)$ is an snc divisor. (By a slight abuse of language, by saying that
a closed subset is an snc divisor we mean that it is the support of an snc divisor.)
If, furthermore, one can always choose such $f$ to be a modification, separable
$l'$-alteration, $l'$-alteration, or alteration, then we say that $X$ is
{\em universally resolvable}, {\em universally separably $l'$-resolvable},
{\em universally $l'$-resolvable}, or {\em universally $\QQ$-resolvable}, respectively.

\begin{remark2}
(i) Due to resolution of singularities in characteristic zero, any qe scheme over
$\Spec(\QQ)$ is universally resolvable for any $l$. This is essentially due to
Hironaka, \cite{Hironaka}, though an additional work was required to treat qe
schemes that are not algebraic in Hironaka's sense, see \cite{Temkin@des1} for
the noetherian case and \cite{Temkin@des2} for the general case.

(ii) It is hoped that all qe schemes admit resolution of singularities
(in particular, are universally resolvable). However, we are, probably,
very far from proving this. Currently, it is known that any qe scheme of
dimension at most two admits resolution of singularities (see \cite{CJS@surfaces}
for a modern treatment). In particular, any qe scheme of dimension at most two is
universally resolvable.

(iii) One can show that any universally $\QQ$-resolvable scheme is qe, but we
prefer not to include this proof here, and will simply add quasi-excellence
assumption to the theorems below.

(iv) On the negative side, we note that there exist regular (hence resolvable)
but not universally $\QQ$-resolvable schemes $X$. They can be constructed analogously
to examples from \refext{I}{1patho-DVR}. For instance, there exists a discretely valued
field $K$ whose completion $\hatK$ contains a finite purely inseparable extension $K'/K$
(e.g. take an element $y\in k((x))$ which is transcendental over $k(x)$ and set
$K=k(x,y^p)\subset  K' = k(x,y) \subset k((x))$ with the induced valuation).
The valued extension $K'/K$ has a defect in the sense that $e_{K'/K}=f_{K'/K}=1$.
In other words, the DVR's $A'$ and $A$ of $K'$ and $K$, have the same residue field
and satisfy $m_{A'}=m_AA'$. Since $A'$ is $A$-flat, it cannot be $A$-finite. On the
other hand, $A'$ is the integral closure of $A$ in $K'$, and we obtain that $A$ is
not qe. In addition, although $X=\Spec(A)$ is regular, any $X$-finite integral scheme
$X'$ with $K'\subseteq k(X')$ possesses a non-finite normalization and hence does not
admit a desingularization. Thus, $X$ is not universally $\QQ$-resolvable.
\end{remark2}

Our main goal will be to show that universal $l'$-resolvability of a qe base scheme
$S$ is inherited by finite type $S$-schemes whose structure morphism $X\to S$ is
maximally dominating (see Theorem \ref{10-l'theor} below, where a more precise
result is formulated). The proof will be by induction on the relative dimension,
and the main work is done when dealing with the case of generically smooth relative curves.

\begin{theorem}\label{10-relcurveth}
Let $S$ be an integral, noetherian, qe scheme with generic point $\eta=\Spec(K)$,
let $f{\colon}X\to S$ be a maximally dominating (\refext{II}{2-maximalement-dominant}) morphism
of finite type, and let $Z\subset X$ be a nowhere dense closed subset. Assume that $S$ is universally
$l'$-resolvable (resp. universally separably $l'$-resolvable), $X_\eta=X\times_S\eta$
is a smooth curve over $K$, and $Z_\eta=Z\times_S\eta$ is étale over $K$. Then there
exist a projective $l'$-alteration (resp. a separable projective $l'$-alteration)
$a{\colon}S'\to S$, a projective modification $b{\colon}X'\to (X\times_SS')^{\rm pr}$,
where $(X\times_SS')^{\rm pr}$ is the proper transform of $X$, i.e. the schematic closure
of $X_\eta\times_SS'$ in $X\times_SS'$,
\begin{equation*}
\xymatrix{X'\ar[r]^b\ar[rrd]_{f'} & (X\times_SS')^{\rm pr}\ar@{^{(}->}[r] & X\times_SS' \ar[d]\ar[r]& X\ar[d]^f\\
& & S' \ar[r]^a & S}
\end{equation*}
and divisors $W'\subset S'$ and $Z'\subset X'$ such that $S'$ and $X'$ are regular,
$W'$ and $Z'$ are snc, the morphism $f'{\colon}X'\to S'$ is pseudo-projective
(\S\ref{10-chowsec}), $(X',Z')\to(S',W')$ is log smooth, and $Z'=c^{-1}(Z)\cup f'^{-1}(W')$,
where $c$ denotes the alteration $X'\to X$.
\end{theorem}
We also note if $f$ is separated (resp. proper) then $f'$ is even quasi-projective
(resp. projective) by Remark \ref{10-chowrem}(ii).
\begin{proof}
It will be convenient to represent $Z$ as $Z_h\cup Z_v$, where the horizontal component
$Z_h$ is the closure of $Z_\eta$ and the vertical component $Z_v$ is the closure of
$Z\setminus Z_h$. The following observation will be used freely: if $a_1{\colon}S_1\to S$
is a (resp. separable) projective $l'$-alteration with an integral $S_1$ and
$b_1{\colon}X_1\to(X\times_SS_1)^{\rm pr}$ is a projective modification, then
it suffices to prove the theorem for $f_1{\colon}X_1\to S_1$ and the preimage
$Z_1\subset X_1$ of $Z$ (note that the generic fiber of $f_1$ is smooth because
it is a base change of that of $f$). So, in such a situation we can freely replace
$f$ by $f_1$, and $Z$ will be updated automatically without mentioning, as a rule.
We will change $S$ and $X$ a few times during the proof. We start with some preliminary steps.

Step 1. {\it We can assume that $f$ is quasi-projective.} By Proposition \ref{10-chowlem},
replacing $X$ with its projective modification we can achieve that $f$ factors through
a local isomorphism $X\to\oX$, where $\oX$ is $S$-projective. Let $X_1\subseteq\oX$ be
the image of $X$. Then the induced morphism $X_1\to S$ is quasi-projective and with
smooth generic fiber. If the theorem holds for $f_1$ and the image $Z_1\subset X_1$
of $Z$, i.e., there exist $a{\colon}S'\to S$ and $b_1{\colon}X'_1\to(X_1\times_SS')^{\rm pr}$
that satisfy all assertions of the theorem, then the theorem also holds for $f$ and $Z$ because
we can keep the same $a$ and take $b=b_1\times_{X_1}X$. This completes the step, and in the
sequel we assume that $f$ is quasi-projective. As we will only use projective modifications
$b_1{\colon}X_1\to(X\times_SS_1)^{\rm pr}$, the quasi-projectivity of $f$ will be preserved
automatically.

Step 2. {\it We can assume that $f$ and $Z_h\to S$ are flat.} Indeed, due to the flattening
theorem of Raynaud-Gruson, see \cite[5.7.9]{Platification@Raynaud-Gruson}, this can be achieved
by replacing $S$ with an appropriate projective modification $S'$, replacing $X$ with the proper
transform, and replacing $Z$ with its preimage. From now on, the proper transforms of $X$ will
coincide with the base changes.

Step 3. {\it Use of the stable modification theorem.} By the stable modification theorem
\cite[1.5 and 1.1]{Temkin@stable} there exist a separable alteration $\oa{\colon}\oS\to S$
with an integral $\oS$ and a projective modification $\oX\to X\times_S\oS$ such that
$(\oX,\oZ_h)$ is a {\em semistable multipointed} $\oS$-curve (see \S\ref{10-relcurvesec}),
where $\oZ_h\subset\oX$ is the horizontal part of the preimage $\oZ$ of $Z$. Enlarging
$\oS$ we can assume that it is integral and normal.
\begin{equation*}
\xymatrix{\oX \ar[rd]_\of\ar[r] & X\times_S\oS \ar[d]\ar[r]& X\ar[d]^f\\
& \oS \ar[r]^\oa & S}
\end{equation*}

Step 4. {\it We can assume that $\oa$ is a separable projective $G$-alteration, where $G$
is an $l$-group.} Since semistable multipointed relative curves are preserved by base changes,
we can just enlarge $\oS$ by replacing it with any separable projective Galois alteration that
factors through $\oS$. Once $\oS\to S$ is Galois, let $\oG$ denote its Galois group and let
$G\subseteq\oG$ be any Sylow $l$-subgroup. Since $\oS\to\oS/G$ is a separable $G$-alteration
and $\oS/G\to S$ is a separable projective $l'$-alteration, we can replace $S$ with $\oS/G$,
replace $X$ with $X\times_S(\oS/G)$, and update $Z$ accordingly, accomplishing the step.

Step 5. {\it The action of $G$ on $X\times_S\oS$ via $\oS$ lifts equivariantly to $\oX$.
In particular, $\of$ becomes $G$-equivariant and $\oX\to X$ becomes a separable projective
$G$-alteration.} This follows from \cite[1.6]{Temkin@stable}.

Step 6. {\it The action of $G$ on $\oX$ is inertia specializing.} Indeed, any covering of
$X$ by separated open subschemes induces a covering of $\oX$ by $G$-equivariant separated
open subschemes. So, it remains to use Remark \ref{10-inspecrem}(i).

Step 7. {\it We can assume that $\oS\to S$ is finite.} By Raynaud-Gruson there exists
a projective modification $S'\to S$ such that the proper transform $\oS'$ of $\oS$ is
flat over $S'$. Let $\eta$ denote the generic point of $S$ and $S'$ and let $\eta'$
denote the generic point of $\oS$ and $\oS'$. Since the morphisms $\oS\times_SS'\to S'$
and $\eta'\to\eta$ are $G$-equivariant and $\oS'$ is the schematic closure of $\eta'$ in
$\oS\times_SS'$, we obtain that the morphism $\oS'\to S'$ is a separable projective $G$-alteration.
Replacing $\oS\to S$ with $\oS'\to S'$, and updating $X$ and $\oX$ accordingly, we achieve that
$\oS\to S$ becomes flat, and hence finite. All conditions of steps 1--6 are preserved with the
only exception that $\oS$ may be non-normal. So, we replace $\oS$ with its normalization and
update $\oX$. This operation preserves the finiteness of $\oS\to S$, so we complete the step.

Step 8. {\it Choice of $W$.} Fix a closed subset $W\subsetneq S$ such that $\oS\to S$ is étale
over $S\setminus W$, $\of(\oZ_v)\subseteq\oW$, where $\oZ_v$ is the vertical part of $\oZ$ and
$\oW=\oa^{-1}(W)$, and $\of$ is smooth over $\oS\setminus\oW$.

Step 9. {\it We can assume that $S$ is regular and $W$ is snc.} Indeed, by our assumptions on
$S$ there exists a projective $l'$-alteration (resp. a separable projective $l'$-alteration)
$a{\colon}S'\to S$ such that $S'$ is regular and $a^{-1}(W)$ is snc. Choose any preimage of
$\eta$ in $S'\times_S\oS$ and let $\oS'$ be the normalization of its closure. Then $\oS'\to S'$
is a separable projective Galois alteration with Galois group $G'\subseteq G$, so we can replace
$S,\oS, G$ and $X$ with $S',\oS', G'$ and $X\times_SS'$, respectively, and update $W$, $\oW$ and
$Z$ accordingly (i.e. replace them with their preimages). Note that step 9 is the only step where
a non-separable alteration of $S$ may occur.

Step 10. {\it The morphism $(\oS,\oW)\to(S,W)$ is Kummer étale.} Indeed, $\oS\to S$ is
an étale $G$-covering outside of $W$, and $\oS$ is the normalization of $S$ in this
covering, so the assertion follows from \refext{IX}{9-FK}.

Consider the $G$-equivariant subscheme $\oT=\oZ\cup\of^{-1}(\oW)$ of $\oX$. The morphism
$(\oX,\oT)\to(\oS,\oW)$ is log smooth by Proposition \ref{10-logsmprop}, hence so is the
composition $(\oX,\oT)\to(S,W)$ and we obtain that $(\oX,\oT)$ is log regular. The group
$G$ acts freely on $\oS\setminus\oW$ and hence also on $\oX\setminus\oT$. Also, its action
on $\oX$ is tame and inertia specializing (step 6), hence we can apply Theorem \ref{10-modif-logsmooth}
to $(\oX,\oT)\to(S,W)$. As a result, we obtain a projective $G$-equivariant modification
$(\oX',\oT')\to(\oX,\oT)$ such that $\oT'$ is the preimage of $\oT$, $G$ acts very tamely
on $(\oX',\oT')$, and $(X',Z')=(\oX'/G,\oT'/G)$ is log smooth over $(S,W)$ (the quotient
exists as a scheme as $f$ is quasi-projective by Step 1 and the morphisms $\oS\to S$ and
$\oX'\to\oX\to X\times_S\oS$ are projective). Clearly, $X'$ is a projective modification
of $X$ and $Z'$ is the union of the preimages of $W$ and $Z$, hence it only remains to
achieve that $X'$ is regular and $Z'$ is snc. For this it suffices to replace $(X',Z')$
with its projective modification $\tilcalF^\rmlog(X',Z')$ introduced in  \refext{VIII}{8-deslogregth}.
\end{proof}

\begin{remark2}
It is natural to compare Theorem \ref{10-relcurveth} and the classical de Jong's result
recalled in \refext{IX}{9-dJ}. The main differences are as follows.

(i) One considers non-proper relative curves in \ref{10-relcurveth}, and this is the only
point that requires to use the stable modification theorem instead of de Jong's result.
The reason is that although the problem easily reduces to the case of a quasi-projective
$f$ (see Step 2), one cannot make $f$ projective, as the compactified generic fiber
$\oX_\eta$ does not have to be smooth (i.e. {\em geometrically} regular) at the added points.

(ii) One uses $l'$-alterations in \ref{10-relcurveth}. This is more restrictive than
in \refext{IX}{9-dJ}, but the price one has to pay is that the obtained log smooth
morphism $(X',Z')\to(S',W')$ does not have to be a nodal curve (e.g. $X'\to S'$ may
have non-reduced fibers). The construction of such $b{\colon}X'\to X$ involves a
quotient by a Sylow subgroup, and is based on Theorem \ref{10-modif-logsmooth}.
(Note also that it seems probable that instead of \ref{10-modif-logsmooth} one
could use a torification argument of Abramovich-de Jong, see \cite[\S1.4.2]{AdJ@Resolution}.)
\end{remark2}

Now, we are going to use Theorem \ref{10-relcurveth} to prove the main result of \S\ref{10-complsec}.

\begin{theorem}\label{10-l'theor}
Let $f{\colon}X\to S$ be a maximally dominating (\refext{II}{2-maximalement-dominant})
morphism of finite type between noetherian qe schemes, let $Z\subset X$ be a nowhere
dense closed subset, and assume that $S$ is universally $l'$-resolvable, then:

(i) $X$ is universally $l'$-resolvable.

(ii) There exist projective $l'$-alterations $a{\colon}S'\to S$ and $b{\colon}X'\to X$
with regular sources, a pseudo-projective (\S\ref{10-chowsec}) morphism $f'{\colon}X'\to S'$
compatible with $f$
\begin{equation*}
\xymatrix{X' \ar[d]^{f'}\ar[r]^b & X \ar[d]^f\\
S' \ar[r]^a & S&}
\end{equation*}
and snc divisors $W'\subset S'$ and $Z'\subset X'$ such that $Z'=b^{-1}(Z)\cup f'^{-1}(W')$
and the morphism $(X',Z')\to(S',W')$ is log smooth.

(iii) If $S=\Spec(k)$, where $k$ is a perfect field, then one can achieve in addition to (ii)
that $a$ is an isomorphism and the alteration $b$ is separable. In particular, $X$ is universally
separably $l'$-resolvable in this case.
\end{theorem}
\begin{proof}
Note that (i) follows from (ii) because any alteration $X_1$ of $X$ is also of finite type
over $S$, so we can apply (ii) to $X_1$ as well. Thus, our aim is to prove (ii) and its
complement (iii). We will view $Z$ both as a closed subset and a reduced closed subscheme.
We start with a few preliminary steps, that reduce the theorem to a special case. We will
tacitly use that if $S_1\to S$ and $X_1\to X$ are projective $l'$-alterations, separable in
case (iii), and $f_1{\colon}X_1\to S_1$ is compatible with $f$, then it suffices to prove the
theorem for $f_1$ and the preimage $Z_1\subset X_1$ of $Z$. So, in such situation we can freely
replace $f$ with $f_1$, and $Z$ will be updated automatically without mentioning, as a rule.

Step 1. {\it We can assume that $X$ and $S$ are integral and normal.} For a noetherian scheme
$Y$ let $\Nor(Y)$ denote the normalization of its reduction. Since $f$ is maximally dominating,
it induces a morphism $\Nor(f){\colon}\Nor(X)\to\Nor(S)$, and replacing $f$ with $\Nor(f)$ we can
assume that $S$ and $X$ are normal. Since we can work separately with the connected components,
we can now assume that $S$ and $X$ are integral.

Step 2. {\it We can assume that $f$ is projective.} By Proposition \ref{10-chowlem} there exists
a projective modification $X_1\to X$ such that the morphism $X_1\to S$ factors into a composition
of a local isomorphism $X_1\to\oX$ and a projective morphism $\of{\colon}\oX\to S$. Replacing $X$
with $X_1$ we can assume that $X$ itself admits a local isomorphism $g{\colon}X\to\oX$ with an
$S$-projective target. Let $\oZ$ be the closure of $g(Z)$. Then it suffices to solve our problem
for $\of$ and $\oZ$, as the corresponding alteration of $\oX$ will induce an alteration of $X$
as required. Thus, replacing $X$ and $Z$ with $\oX$ and $\oZ$, we can assume that $f$ is projective.

Step 3. {\it It suffices to find $f'$ which satisfies all assertions of the theorem except the
formula for $Z'$, while the latter is weakened as $b^{-1}(Z)\cup f'^{-1}(W')\subseteq Z'$.} Given
such an $f'$ note that $Z''=b^{-1}(Z)\cup f'^{-1}(W')$ is a subdivisor of $Z'$, hence it is an snc
divisor too. We claim that $X',Z'',S',W'$ satisfy all assertions of the theorem, and the only thing
one has to check is that the morphism $(X',Z'')\to(S',W')$ is log smooth. The latter follows from
Lemma \ref{10-logsmlem} whose proof will be given below.

Step 4. {\it In the situation of (iii) we can assume that the field $k$ is infinite.} Assume that
$S=\Spec(k)$ where $k$ is a finite field and fix an infinite algebraic $l'$-prime extension $\ok/k$
(i.e. it does not contain the extension of $k$ of degree $l$). We claim that it suffices to prove
(ii) and (iii) for $\oS=\Spec(\ok)$ and the base changes $\oX=X\times_S\oS$ and $\oZ=Z\times_S\oS$.
Indeed, assume that $a{\colon}\oX'\to\oX$ is a separable $l'$-alteration with a regular source and
such that $\oZ'=a^{-1}(\oZ)$ is an snc divisor (obviously, we can take $\oS'=\oS$ and $\oW'=\emptyset$).
Since $\oS=\lim_i S_i$ where $S_i=\Spec(k_i)$ and $k_i/k$ run over finite subextensions of $\ok/k$,
\cite[8.8.2(ii)]{EGAIV-3} implies that there exists $i$ and a finite type morphism
$X'_i\to X_i=X\times_SS_i$ such that $\oX'\iso X'_i\times_{S_i}\oS$. For any finite
subextension $k_i\subseteq k_j\subset\ok$ set $X'_j=X'_i\times_{S_i}S_j$ and $X_j=X_i\times_{S_i}S_j$,
and let $Z'_j\subset X'_j$ be the preimage of $Z$. Then it follows easily from \cite[8.10.5]{EGAIV-3}
and \cite[17.7.8]{EGAIV-4} that $X'_j\to X_j$ is an $l'$-alteration and $X'_j\to S_j$ is smooth for
large enough $k_j$. In the same manner one achieves that $Z'_j$ is an snc divisor. Now, it is obvious
that $(X'_j,Z'_j)\to(S,\emptyset)$ is log smooth and $X'_j\to X_j\to X$ is an $l'$-alteration.

Now we are in a position to prove the theorem. We will use induction on $d=\trdeg(k(X)/k(S))$, with
the case of $d=0$ being obvious. Assume that $d\ge 1$ and the theorem holds for smaller values of $d$.

Step 5. {\it Factorizing $f$ through a relative curve.} After replacing $X$ by a projective modification
we can factor $f$ through an integral scheme $Y$ such that $g{\colon}X\to Y$ is maximally dominating,
$h{\colon}Y\to S$ is projective, and $\trdeg(k(X)/k(Y))=1$. Indeed, one can obviously construct such
a rational map $h'{\colon}X\dasharrow Y$ even without modifying $X$ (i.e. $h'$ is well defined on a
non-empty open subscheme $U\subseteq X$). Let $X'$ be the schematic image of the morphism $U\into X\times Y$.
Then $X'\to X$ is a projective modification (an isomorphism over $U$), and the morphism $X'\to S$
factors through $Y$.

Let $\eta=\Spec(k(Y))$ denote the generic point of $Y$, $X_\eta=X\times_{Y}\eta$ and
$Z_\eta=Z\times_Y\eta$. If $k$ is an infinite perfect field and $S=\Spec(k)$ then
a Bertini type argument from the end of proof of \cite[4.11]{Alterations@de_Jong}
implies that we can choose $Y$ so that $X_\eta$ and $Z_\eta$ are smooth over $\eta$.
(One assumes that $k$ is algebraically closed in \cite{Alterations@de_Jong}, but to
use Bertini's theorem one only needs that $k$ is perfect and infinite.) In the general
case, we can achieve that $X_\eta$ and $Z_\eta$ are smooth at cost of a purely inseparable
alteration of $Y$. Indeed, pick up any finite purely inseparable extension $K/k(Y)$ such
that $Z_K=\Nor(Z\times_\eta\Spec(K))$ (i.e. just the reduction) and $X_K=\Nor(X\times_\eta\Spec(K))$
are smooth, extend $K$ to a projective alteration $Y'\to Y$, and replace $Y$ and $X$ with $Y'$
and the schematic closure of $X_K$ in $X\times_YY'$, respectively.

Step 6. {\it Use of Theorem \ref{10-relcurveth}.} So far, we have constructed the right column
of the following diagram
$$
\xymatrix{(X',Z') \ar[rrdd]_{f'}\ar[rr]^{\tilcalF^\rmlog(L,M_L)} & &
(L,M_L) \ar[r]\ar[d]^{g'}\ar@{}[rd]|\square & (X'',Z'') \ar[d]^{g''}\ar[r]^{b''} & X \ar[d]^g\ar@/^1pc/[dd]^f\\
 & & (Y',V')\ar[r]^{c'}\ar[d]^{h'} & (Y'',V'')\ar[r]^{c''}\ar[d]^{h''} & Y\ar[d]^h\\
 & & (S',W')\ar[r]^a & S\ar@{=}[r] & S
}
$$
By Theorem \ref{10-relcurveth} there exists a projective $l'$-alteration $c''{\colon}Y''\to Y$
with regular source, a projective modification $X''\to(X\times_YY'')^{\rm pr}$ with regular
source, a projective morphism $g''{\colon}X''\to Y''$ compatible with $g$, and snc divisors
$V''\subset Y''$ and $Z''\subset X''$ such that $(X'',Z'')\to(Y'',V'')$ is log smooth and
$b''^{-1}(Z)\subseteq Z''$. In case (iii), $Y$ is universally separably $l'$-resolvable by
the induction assumption, hence we can take $c''$ to be separable, and then $b''{\colon}X''\to X$
is also separable. In addition, by the induction assumption applied to $h''{\colon}Y''\to S'$ and
$V''\subset Y''$ there exist projective $l'$-alterations $a{\colon}S'\to S$ and
$c'{\colon}Y'\to Y''$ with regular sources and snc divisors $W'\subset S'$ and
$V'\subset Y'$ and a projective morphism $h'{\colon}Y'\to S'$ compatible with $h''$
such that $(Y',V')\to(S',W')$ is log smooth, $c'^{-1}(V'')\subseteq V'$, and $c'$ is
separable if the assumption of (iii) is satisfied.

Set $(L,M_L)=(X'',Z'')\times^{\rm fs}_{(Y'',V'')}(Y',V')$, where the product is
taking place in the category of fs log schemes. To simplify notation we will write
$P^\sat$ instead of $(P^{\rm int})^\sat$ for monoids and log schemes. Recall that
$(L,M_L)=(F,M_F)^\sat$, where $(F,M_F)$ is the usual log fibered product, and
$F=X''\times_{Y''}Y'$ by \cite[1.6]{LSFI@Kato}. Furthermore, we have local Zariski
charts for $c'$ and $g''$ modeled, say, on $P_i\to P'_i$ and $P_i\to Q_i$. Hence
$(F,M_F)$ is a Zariski log scheme with charts modeled on $R_i=P'_i\oplus_{P_i}Q_i$,
and $(L,M_L)$ is a Zariski log scheme with charts modeled on $R_i^\sat$. Furthermore,
the saturation morphism $L\to F$ is finite hence the composition $L\to F\to X''$ is
projective. The morphism $g'{\colon}(L,M_L)\to (Y',V')$ is a saturated base change of
the log smooth morphism $g''{\colon}(X'',Z'')\to(Y'',V'')$, hence it is log smooth.
As $(Y',V')$ is log regular, $(L,M_L)$ is also log regular. Applying to $(L,M_L)$ the
saturated monoidal desingularization functor $\tilcalF^\rmlog$ from \refext{VIII}{8-deslogregth}
we obtain a log regular Zariski log scheme $(X',Z')$ with a regular $X'$. Then $Z'$ is a normal
crossings divisor, which is even an snc divisor since the log structure is Zariski.

We claim that $(X',Z')$ and $(S',W')$ are as asserted by the theorem except of the
weakening dealt with in Step 3. Indeed, the morphism $(X',Z')\to(S',W')$ is log
smooth because it is the composition $(X',Z')\to (L,M_L)\to(Y',V')\to(S',W')$ of
log smooth morphisms. The preimage of $Z$ in $X''$ is contained in $Z''$, which is
the non-triviality locus of the log structure of $(X'',Z'')$, hence its preimage in
$X'$ is also contained in the non-triviality locus of the log structure of $(X',Z')$,
which is $Z'$. Clearly, $Z'$ also contains the preimage of $W'$. By the construction,
$S'\to S$ is a projective $l'$-alteration, and it remains to check that $X'\to X$ is
also a projective $l'$-alteration. Since $\tilcalF^\rmlog(L,M_L)$ is a saturated log
blow up tower and $(L,M_L)$ is log regular, the underlying morphism of schemes $X'\to L$
is a projective modification by \refext{VIII}{8-logregblowlem}(i). The projective morphism
$L\to X''$ is an $l'$-alteration because generically (where the log structures are trivial)
it is a base change of the projective $l'$-alteration $Y'\to Y''$. And $X''\to X$ is
a projective $l'$-alteration by the construction. It remains to recall that in the
situation of (iii) the alterations $c'{\colon}Y'\to Y''$ and $b''{\colon}X''\to X$
are separable, hence so are $(L,M_L)\to X''$ and the total composition $X'\to X$.
\end{proof}

\begin{remark2}
The only place where inseparable alterations are used is the argument at step 5,
where we had to choose an inseparable extension $K/k(Y)$ when $X_\eta$ or $Z_\eta$
is not geometrically regular.
\end{remark2}

\begin{lemma2}\label{10-logsmlem}
Assume that $S$ and $X$ are regular schemes, $W\subset S$ and $Z\subset X$ are snc
divisors, and $f:X\to S$ is a morphism such that $f^{-1}(W)\subseteq Z$ and the
induced morphism of log schemes $h{\colon}(X,Z)\to(S,W)$ is log smooth. Then for
any intermediate divisor $f^{-1}(W)\subseteq Z'\subseteq Z$ the morphism $h'{\colon}(X,Z')\to(S,W)$
is log smooth.
\end{lemma2}
\begin{proof}
We can work locally at a geometric point $\ox\to X$. Let $x\in X$ and $s\in S$
be the images of $\ox$, and let $q_1\. q_r\in\calO_{S,s}$ define the irreducible
components of $W$ through $s$. Set $Q=\oplus_{i=1}^r q_i^\NN$. Shrinking $S$ we
obtain a chart $c{\colon}S\to\Spec(\ZZ[Q])$ of $(S,W)$. By Proposition \ref{10-local-structure}
applied to $c$, $h$, and $G=1$, after localizing $X$ along $\ox$ one can find an fs
chart of $h$ consisting of $c$, $X\to\Spec(\ZZ[P])$, and $\phi{\colon}Q\to P$ such
that the morphism $X\to S\times_{\Spec(\ZZ[Q])}\Spec(\ZZ[P])$ is smooth, $P^*$ is
torsion free, $\phi$ is injective, and the torsion of $\Coker(\phi^\gp)$ is annihilated
by an integer $n$ invertible on $S$.

Let $p_1\. p_t\in\calO_{X,x}$ define the irreducible components of $Z$ through $x$.
Our next aim is to adjust the chart similarly to \ref{10-chart-diagram}(vi) to achieve
that $\oP\iso N=\oplus_{j=1}^t p_j^\NN$. Note that $\oM_{X,x}\iso N$, where
$M_X\into\calO_X$ is the log structure of $(X,Z)$. The homomorphism $\psi{\colon}P\to M_{X,x}$
factors through the fs monoid $R=P[T^{-1}]$ where $T=\psi^{-1}(M^*_{X,x})$. Clearly,
$\R^*$ is torsion free, $\oR\iso N$, and shrinking $X$ around $x$ we obtain a chart
$X\to\Spec(\ZZ[R])$. Since $\Spec(\ZZ[R])$ is open in $\Spec(\ZZ[P])$ the morphism
$X\to S\times_{\Spec(\ZZ[Q])}\Spec(\ZZ[R])$ is smooth. So, we can replace $P$ with
$R$ achieving that $\oP\iso N$, and hence $P\iso\NN^t\oplus\ZZ^u$.

Without restriction of generality $Z'$ is defined by the vanishing of $\prod_{j=1}^{t'}p_j$
for $t'\le t$. Since $f^{-1}(W)\subseteq Z'$, the image of $Q$ in $\oP$ is contained in
$\oP'=\oplus_{j=1}^{t'}p_j^\NN$. Hence $\phi$ factors through a homomorphism
$\phi'{\colon}Q\to P'=\oP'\oplus P^*$, and we obtain a chart of $h'$ consisting
of $c$, $X\to\Spec(\ZZ[P'])$, and $\phi'$. By \cite[3.5,3.6]{Toric@Kato},
to prove that $h'$ is log smooth it remains to observe that $\phi'$ is injective,
the torsion of $\Coker(\phi'^\gp)$ is annihilated by $n$, and the morphism
$$X\to S\times_{\Spec(\ZZ[Q])}\Spec(\ZZ[P])\to S\times_{\Spec(\ZZ[Q])}\Spec(\ZZ[P'])$$
is smooth because $\Spec(\ZZ[P])\to\Spec(\ZZ[P'])$ is so.
\end{proof}

\subsubsection{Comparison with Theorems \ref{10-de Jong1} and \ref{10-de Jong2}}
Theorems \ref{10-de Jong1} and \ref{10-de Jong2} follow by applying Theorem \ref{10-l'theor}(ii)
to $X\to S$ and $Z$ (where one takes $S=\Spec(k)$ in \ref{10-de Jong1}). Indeed,
the main part of the proofs of \ref{10-de Jong1} and \ref{10-de Jong2} was to
construct $l'$-alterations $X'\to X$ and $S'\to S$ with regular sources, snc
divisors $Z'\subset X'$ and $W'\subset S'$, and a log smooth morphism
$f'{\colon}(X',Z')\to(S',W')$ compatible with $f$. Then, in the last paragraphs
of both proofs, Proposition \ref{10-local-structure} was used to obtain a more
detailed description of $X'$, $Z'$, and $f'$. In particular, for a zero-dimensional
base this amounted to saying that $X'$ is $S$-smooth and $Z'$ is relatively snc over
$S$, and for a one-dimensional base this amounted to the conditions (i) and (ii)
of \ref{10-de Jong2}.

Conversely, Theorem \ref{10-de Jong1} (resp. \ref{10-de Jong2}) implies assertion (ii)
of Theorem \ref{10-l'theor} under the assumptions of \ref{10-de Jong1} (resp. \ref{10-de Jong2})
on $X$ and $S$. Moreover, the non-separated Chow's lemma could be used in their proofs as well,
so the separatedness assumption there could be easily removed. In such case,
Theorems \ref{10-de Jong1} and \ref{10-de Jong2} would simply become the low
dimensional (with respect to $S$) cases of Theorem  \ref{10-l'theor}(ii) plus
an explicit local description of the log smooth morphism $f'$. The strengthening
\ref{10-l'theor}(iii), however, was not achieved in \ref{10-de Jong1}, and required
a different proof of the whole theorem.

\subsection{Saturation}\label{10-satursec} In Theorems \ref{10-relcurveth} and \ref{10-l'theor}
we resolve certain morphisms $f{\colon}X\to S$ with divisors $Z\subset X$ by log smooth morphisms
$f'{\colon}(X',Z')\to(S',W')$. However, as we insisted to use only $l'$-alterations and to obtain
regular $X'$ and snc $Z'$, we had to compromise a little on the "quality" of $f'$. For example,
our $f'$ may have non-reduced fibers. Due to de Jong's theorem, if the relative dimension is one,
then one can make $f'$ a nodal curve. We will see that a similar improvement of $f'$ is possible
in general if one uses arbitrary alterations and allows non-regular $X'$. The procedure reduces
to saturating $f'$ and is essentially due to Tsuji and Illusie-Kato-Nakayama.

\subsubsection{Saturated morphisms} Recall that a homomorphism $P\to Q$ of fs (resp. fine)
monoids is {\em saturated} (resp. {\em integral}) if for any homomorphism $P\to P'$ with
fs (resp. fine) target the pushout $Q\oplus_{P}P'$ is fs (resp. fine). A morphism of fs (resp. fine)
log schemes $f{\colon}(Y,M_Y)\to(X,M_X)$ is {\em saturated} (resp. {\em integral}) if so are
the homomorphisms $\oM_{X,f(y)}\to\oM_{Y,y}$.

\begin{remark2}\label{10-satrem}
(i) Integral morphism were introduced already by Kato in \cite[\S4]{LSFI@Kato}. Kato also
introduced the notion of saturated morphisms, which was first seriously studied by Tsuji
in \cite{Saturated@Tsuji}. Actually, one can define saturated morphisms for arbitrary fine
log schemes, but the definition is more involved than we use. For fs log schemes our definition
coincides with the usual one due to \cite[II 2.13(2)]{Saturated@Tsuji}.

(ii) The following two basic properties of saturated morphisms follow from the definition:
(a) a composition of saturated morphisms between fs log schemes is saturated,
(b) if $f{\colon}Y\to X$ is a saturated morphism between fs log schemes, then,
for any morphism of fs log schemes $Y'\to Y$, the base change $f'{\colon} Y'\to X'$ of $f$
in the category of log schemes is a saturated morphism of fs log schemes. (Also, it is proved
in \cite[II 2.11]{Saturated@Tsuji} that analogous properties hold for saturated morphisms between
arbitrary integral log schemes.)

(iii) Let $f{\colon}Y\to X$ be a morphism of fs log schemes. It is shown in \cite[II 3.5]{Saturated@Tsuji}
that if $f$ can be modeled on charts corresponding to saturated homomorphisms of fs monoids $P_i\to Q_i$
then $f$ is saturated. Let us remark that the converse is also true: if $f$ is saturated then it can be
modeled on charts corresponding to $P_i\to Q_i$ as above.
\end{remark2}

\subsubsection{Integrality and saturatedness for log smooth morphisms}
 We recall the following result that relates the notions of integral and saturated
 morphisms to certain properties of the underlying morphisms of schemes.

\begin{proposition2}\label{10-saturprop1}
Let $f{\colon}(Y,M_Y)\to(X,M_X)$ be a log smooth morphism between fs log schemes.

(i) If $f$ is integral then $Y\to X$ is flat.

(ii) Assume that $f$ is integral. Then $f$ saturated if and only if $Y\to X$ has reduced fibers.
\end{proposition2}
\begin{proof}
The first claim is proved in \cite[4.5]{LSFI@Kato} and the second one is proved in \cite[II 4.2]{Saturated@Tsuji}.
\end{proof}

Sometimes one can also go in the opposite direction: from flatness to integrality.

\begin{proposition2}\label{10-saturprop2}
Let $f{\colon}(Y,M_Y)\to(X,M_X)$ be a log smooth morphism between fs log schemes and
 assume that the morphism $Y\to X$ is flat, then $f$ is integral.
\end{proposition2}
\begin{proof}
It suffices to show that if $\oy\to Y$ is a geometric point and $\ox=f(\oy)$ then
the homomorphism $\ophi{\colon}\oM_{X,\ox}\to\oM_{Y,\oy}$ is integral. By Proposition
\ref{10-local-structure} and the argument in \ref{10-preliminary}(vi), localizing $X$
and $Y$ along these points we can assume that $f$ is modeled on a chart
$Y_0=\Spec(\ZZ[P])\to X_0=\Spec(\ZZ[Q])$ corresponding to a homomorphism
$\phi{\colon}Q\to P$ so that the morphism $g{\colon}Y\to Z=X\times_{X_0}Y_0$
is smooth (in particular, flat), and $\phi$ has the following properties:
$Q =\oM_{X,\ox}$, $P$ is fs, $P^*$ is torsion free, the composition
$Q\to P \to\oP = P/P^*$ coincides with $\ophi$, the kernel of $\phi^{\gp}$ is finite,
killed by an integer invertible at $x$, as well as the torsion part of its cokernel
(but we will not need these last two properties). Since $Q$ is sharp and saturated,
$Q^{\gp}$ is torsion free, so $\phi$ is injective. We claim that $\ophi$ is is integral
if and only if $\phi$ is integral. To see this note that if $Q\to R$ is a homomorphism
of monoids, then $R\oplus_Q\oP$ is isomorphic to the quotient of $R\oplus_Q P$ by the
image of $P^*$, and hence either both pushouts are integral or neither of them is integral.
Thus, we need only prove that $\phi$ is integral.

Note that the morphism $h{\colon}Z\to X$ is flat at the (Zariski) image $z\in Z$ of $\oy$
because $f$ and $g$ are flat. Set $x=h(z)$ and $k=k(x)$, then the fiber
$h_x{\colon}\Spec(k[P])\to\Spec(k[Q])$ of $h$ over $x$ is flat at $z$.
In other words, if $I\subset k[P]$ is the ideal corresponding to $z$ then
the homomorphism $k[Q]\to k[P]_I$ is flat. The preimage of $m_y$ under
$\ZZ[P]\to\calO_{Y,y}$ contains the set $m_P=P\setminus P^*$, hence $m_P\subset I$
and we obtain that $I$ contains $J=k[m_P]$. Note that the ideal $J$ is prime as
$k[P]/J\iso k[P^*]$ is a domain due to $P^*$ being torsion free. Thus, the
localization $k[P]_J$ makes sense, and we obtain a flat homomorphism $\psi{\colon}k[Q]\to k[P]_J$.

It is proved in \cite[4.1]{LSFI@Kato}, that if the homomorphisms $K[\phi]{\colon}K[Q]\to K[P]$
are flat for any field $K$ then $\phi$ is integral. The proof consists of two parts. First one
checks that $\phi$ is injective, which is automatic in our case. This is the only argument in
loc.cit. where a play with different fields is needed. We claim that the second part of the
proof of the implication (iii)$\implies$(v) in \cite[4.1]{LSFI@Kato} works fine with a single
field $k$, and, moreover, it suffices to only use that $k[Q]\to k[P]_J$ is flat. Let us indicate
how the argument in loc.cit. should be adjusted.

Assume that, as in the proof of \cite[4.1]{LSFI@Kato}, we are given $a_1,a_2\in Q$ and
$b_1,b_2\in P$ such that $\phi(a_1)b_1=\phi(a_2)b_2$. Let $S$ be the kernel of the
homomorphism of $k[Q]$-modules $k[Q]\oplus k[Q]\to k[Q]$ given by $(x,y)\mapsto a_1x-a_2y$.
By the flatness, the kernel of $k[P]_J\oplus k[P]_J\to k[P]_J$, $(x,y)\mapsto\phi(a_1)x-\phi(a_2)y$
is generated by the image of $S$. Hence there exist representations
$b_1=\sum_{i=1}^r\phi(c_i)\frac{f_i}{s}$ and $b_2=\sum_{i=1}^r\phi(d_i)\frac{f_i}{s}$
with $c_i,d_i\in k[Q]$, $f_i\in k[P]$, $s\in k[P]\setminus J$, and $a_1c_i=a_2d_i$.
Moreover, multiplying $s$ and $f_i$'s by an appropriate unit $u\in P^*$ we can assume
that $s=1+s'$ for $s'\in{\rm Span}_k(P\setminus\{1\})$. Then
$b_1+\sum_{1\le\alpha\le m}\lambda_\alpha t_\alpha=\sum_{1\le i \le r}\phi(c_i)f_i$,
with $\lambda_\alpha\in k^*$, and the $t_\alpha\in P$ pairwise distinct and distinct
from $b_1$, so we see that there exist $a_3\in Q$, $b\in P$, and $1\le i\le r$, such
that $a_3$ appears in $c_i$, $b$ appears in $f_i$, and $b_1=\phi(a_3)b$. The remaining
argument copies that of the loc.cit. verbatim, and one obtains in the end that $\phi$
satisfies the condition (v) from \cite[4.1]{LSFI@Kato}. Thus, $\phi$ is integral and
we are done.
\end{proof}

Before going further, let us discuss an incarnation of saturated morphisms in (more classical)
toroidal geometry.

\begin{remark2}
In toroidal geometry an analog of saturated morphisms was introduced by Abramovich and Karu
in \cite{AK@Toroidal}. In the language of log schemes toroidal morphisms can be interpreted
as log smooth morphisms $f{\colon}(X,Z)\to(S,W)$ between log regular schemes (with the toroidal
structure given by the triviality loci of the log structures). If  $f$ is a toroidal morphism
as above then Abramovich-Karu called it weakly semistable when the following conditions hold:
$S$ is regular, $f$ is locally equidimensional, and the fibers of $f$ are reduced. Furthermore,
they remarked that the equidimensionality condition is equivalent to flatness of $f$ whenever
$S$ is regular, see \cite[4.6]{AK@Toroidal}. Thus, the weak semistability condition is nothing
else but saturatedness of $f$ and regularity of the target. In particular, saturated log smooth
morphisms between log regular log schemes may be viewed as the generalization of weakly semistable
morphisms to the case of an arbitrary log regular (or toroidal) base.
\end{remark2}

Now, we are going to prove the main result of \S\ref{10-satursec}.

\begin{theorem}\label{10-saturth}
Assume that $f{\colon}(X,Z)\to(S,W)$ is a log smooth morphism such that $(S,W)$ is log regular
and $S$ is universally $\QQ$-resolvable (\S\ref{10-resolvsec}). Then there exists an alteration
$h{\colon}S'\to S$ such that $S'$ is regular, $W'=g^{-1}(W)$ is an snc divisor, and the fs base
change $f'{\colon}(X',Z')\to(S',W')$ is a saturated morphism.
\end{theorem}
Recall that $(X',Z')=(X,Z)\times_{(S,W)}^{\rm fs}(S',W')$ and $f'$ is log smooth because the
saturation morphism is log smooth.
\begin{proof}
By \refext{VIII}{8-deslogregth}, applying to $(S,W)$ an appropriate saturated log blow up
tower and replacing $(X,Z)$ with the fs base change we can achieve that $S$ is regular and
$W$ is normal crossings. By an additional sequence of log blow ups we can also make $W$ snc
(see \refext{VIII}{8-step2}), so $(S,W)$ becomes a Zariski log scheme. Now, we can étale-locally
cover $f$ by charts $f_i{\colon}(X_i,Z_i)\to(S_i,W_i)$ modeled on $P_i\to Q_i$ such that $S_i$ are
open subschemes in $S$. By \cite[A.4.4, A.4.3]{IKN@Quasi-unipotent}, for each $i$ there exists
a morphism $h_i{\colon}(S'_i,W'_i)\to(S_i,W_i)$ such that $h_i$ is a composition of a Kummer
morphism and a log blow up, and the fs base change of $f_i$ is saturated. (Although the proof
in loc.cit. is written in the context of log analytic spaces, it translates to our situation
almost verbatim. The only changes are that we have to distinguish étale and Zariski topology
on the base (in order to construct log blow ups), and $h_i$ does not have to be log étale as
there might be inseparable Kummer morphisms.)

Note that $W'_i=h_i^{-1}(W_i)$. In addition, $S'_i\to S_i$ is a projective alteration by
\refext{VIII}{8-logregblowlem}. Extend each $h_i$ to a projective alteration $g_i{\colon}T_i\to S$,
and let $h{\colon}S'\to S$ be a projective alteration that factors through each $T_i$.
By the universal $\QQ$-resolvability assumption we can enlarge $S'$ so that it becomes
regular and $Z'=h^{-1}(Z)$ becomes snc. We claim that $h$ is as claimed. It suffices to
check that the fs base change of each morphism $f'_i{\colon}(X_i,Z_i)\to(S,W)$ is saturated.
However, already the fs base change of $f'_i$ to $(T_i,g_i^{-1}(W))$ is saturated by the construction,
hence so is its further base change to $(S',W')$.
\end{proof}

\begin{remark2}
Our proof is an easy consequence of \cite[A.4.4 and A.4.3]{IKN@Quasi-unipotent}. The first cited
result shows that (locally) any log smooth morphism can be made exact by an appropriate log blow
up of the base. This result is somewhat analogous to the flattening theorem of Raynaud-Gruson.
The second cited result shows that by a Kummer extension of the base one can (locally) saturate
an exact log smooth morphism. It is somewhat analogous to the reduced fiber theorem of Bosch-L\"utkebohmert-Raynaud
(\cite{BLR@Reduced_Fiber}) which implies that if $f{\colon}Y\to X$ is a finite type morphism between
reduced noetherian schemes then there exists an alteration $X'\to X$ such that the normalized base
change $f'{\colon}\Nor(Y\times_XX')\to X'$ has reduced fibers. Although the proof of the latter is
far more difficult.
\end{remark2}

\subsection{Characteristic zero case}\label{10-char0sec} Theorem \ref{10-l'theor} can be substantially
strengthened when $S$ is of characteristic zero, i.e., the morphism $S\to\Spec(\ZZ)$ factors through $\Spec(\QQ)$.

\begin{theorem}\label{10-char0theor}
Assume that $S$ is a reduced, noetherian, qe scheme of characteristic zero, $f{\colon}X\to S$
is a maximally dominating morphism of finite type with reduced source, and $Z\subset X$ is a
nowhere dense closed subset. Then there exist projective modifications $a{\colon}S'\to S$ and
$b{\colon}X'\to X$ with regular sources, a pseudo-projective morphism $f'{\colon}X'\to S'$
compatible with $f$, and snc divisors $W'\subset S'$ and $Z'\subset X'$ such that
$Z'=b^{-1}(Z)\cup f'^{-1}(W')$ and the morphism $(X',Z')\to(S',W')$ is log smooth.
\end{theorem}
\begin{proof}
The proof is very close to the proof of Theorem \ref{10-l'theor}, so we will just say which changes
in that proof should be made. First, we note that any $S$-scheme $Y$ of finite type is noetherian and qe.
Thus, if $Y$ is reduced and $T\subset Y$ is a nowhere dense closed subset then the pair $(Y,T)$
can be desingularized by \cite{Temkin@des1} in the following sense: there exists a projective
modification $h{\colon}Y'\to Y$ with regular source and such that $h^{-1}(T)$ is an snc divisor.
This result replaces the $l'$-resolvability assumption in Theorem \ref{10-l'theor}, and it allows
to apply the proof of that theorem to our situation with the only changes that one always uses
projective modifications instead of projective $l'$-alterations, and Theorem \ref{10-relcurveth}
is replaced with Lemma \ref{10-relcurvelem} below. (Note that Lemma \ref{10-relcurvelem}
is weaker than Theorem \ref{10-char0theor}, while Theorem \ref{10-relcurveth} does not follow
from Theorem \ref{10-l'theor}.)
\end{proof}

\begin{lemma2}\label{10-relcurvelem}
Let $S$ be an integral, noetherian, qe scheme with generic point $\eta=\Spec(K)$,
let $f{\colon}X\to S$ be a maximally dominating morphism of finite type, and let
$Z\subset X$ be a nowhere dense closed subset. Assume that $X_\eta=X\times_S\eta$
is a smooth curve over $K$, and $Z_\eta=Z\times_S\eta$ is étale over $K$. Then
there exist projective modifications $a{\colon}S'\to S$ and $b{\colon}X'\to X$
with regular sources, a pseudo-projective morphism $f'{\colon}X'\to S'$ compatible
with $f$ and snc divisors $W'\subset S'$ and $Z'\subset X'$ such that
$Z'=b^{-1}(Z)\cup f'^{-1}(W')$ and the morphism $(X',Z')\to(S',W')$ is log smooth.
\end{lemma2}
\begin{proof}
The proof copies the proof of Theorem \ref{10-relcurveth} with the only difference
that instead of an $l$-Sylow subgroup $G\subseteq\oG$ one simply takes $G=\oG$.
The latter is possible because the schemes are of characteristic zero and hence
any action of $\oG$ is tame.
\end{proof}

Combining Theorem \ref{10-char0theor} and \ref{10-saturth} we obtain the following
weak semistable reduction theorem.

\begin{theorem}\label{10-weaksemistable}
Assume that $S$ is a reduced, noetherian, qe scheme of characteristic zero,
$f{\colon}X\to S$ is a maximally dominating morphism of finite type with
reduced source, and $Z\subset X$ is a nowhere dense closed subset. Then
there exists an alteration $S'\to S$, a modification $X'\to X\times_SS'$
of the proper transform of $X$, a pseudo-projective morphism $f'{\colon}X'\to S'$
compatible with $f$, and divisors $W'\subset S'$ and $Z'\subset X'$ such that $S'$
is regular, $W'$ is snc, $Z'=b^{-1}(Z)\cup f'^{-1}(W')$, and the morphism $(X',Z')\to(S',W')$
is log smooth and saturated (i.e. $X'\to S'$ is weakly semistable).
\end{theorem}

\begin{remark2}
(i) In the case when $X$ and $S$ are integral proper varieties over an algebraically
closed field $k$ of characteristic zero, this theorem becomes the weak semistable
reduction theorem of Abramovich-Karu. Our proof has many common lines with their
arguments. In particular, the first step of their proof was to make $f$ toroidal,
and it was based on de Jong's theorem. (Note also that in a recent work \cite{ADK@Toroidal}
of Abramovich-Denef-Karu, the toroidalization theorem was extended to separated schemes
of finite type over an arbitrary ground field of characteristic zero.)
Our Theorem \ref{10-char0theor} can be viewed as a generalization
of the toroidalization theorem of Abramovich-Karu.

(ii) The second stage in the proof of the weak semistable reduction theorem of Abramovich-Karu
(the combinatorial stage) is analogous to Theorem \ref{10-saturth}. It obtains as an input
a toroidal morphism $f{\colon}(X,Z)\to(S,W)$ between proper varieties of characteristic zero
and outputs an alteration $h{\colon}S'\to S$ such that $S'$ is regular, $W'=h^{-1}(W)$ is snc,
and the saturated base change of $f$ is weakly semistable. The proof is similar to the arguments
used in the proofs of \cite[A.4.4 and A.4.3]{IKN@Quasi-unipotent}. First, one constructs a toroidal
blow up of the base that makes the fibers equidimensional (i.e. makes the log morphism integral),
and then an appropriate normalized finite base change is used to make the fibers reduced.
\end{remark2}

\ifx\danslelivre\undefined
\bibliography{../fichiers-configuration-latex/bibliographie}
\bibliographystyle{../fichiers-configuration-latex/smfalpha}
\end{document}
\fi
